% =====================================================================================================
% ch08_pdf.tex  —  第 8 章  严格核心 III：部分子分布与碎裂函数的算符定义
% writer W09 (spec v2.3.1; LEAD_DECISIONS LA-1…LA-16). Scripts and logs: writing/work/W09/ (v1–v10).
% =====================================================================================================
\chapter{严格核心 III：部分子分布与碎裂函数的算符定义}\label{ch:pdf}

本章把部分子分布函数（parton distribution function，PDF）与碎裂函数（fragmentation function，FF）定义为规范不变、经过重整化的算符矩阵元，并在任何因子化定理之前证明它们的基本性质。这一次序本身是严格性的要求：若把 PDF 定义为“DIS 中乘在系数函数前面的量”，普适性就成了定义而不是定理。各结论的严格性等级见下面的方框；推广到强子态的结论另需\cref{asm:intro-wightman}、\cref{asm:intro-composite}。矩关系的全阶证明见\cref{app:appL}，所用的矩问题引理见\cref{app:appK}。

\begin{tcolorbox}[enhanced,breakable,title=本章证明了什么／没有证明什么,fonttitle=\bfseries,
  colback=gray!3!white,colframe=gray!55!black,boxrule=0.6pt,arc=2pt,
  left=5pt,right=5pt,top=3pt,bottom=3pt,before skip=6pt,after skip=6pt]
\textbf{证明了：}\cref{prop:pdf-invariance}、\cref{lem:pdf-wilson-moments}（\tagR，$R_{\mathrm{def}}$）；\cref{thm:pdf-support-pt}、\cref{thm:pdf-moments-local}、\cref{thm:pdf-dglap}、\cref{prop:pdf-lo-kernels}、\cref{prop:pdf-pqq-operator}、\cref{prop:pdf-partonic-dr}、\cref{prop:pdf-local-anomalous}、\cref{prop:pdf-sum-rules}（\tagR，$R_{\mathrm{PT}}$；强子版本需\cref{asm:intro-composite}）；\cref{thm:pdf-dis-xspace}（\tagR，$R_{\mathrm{OPE}}$）；\cref{prop:pdf-timelike-dglap}的单圈内容(b)（\tagR，$R_{\mathrm{PT}}$）。\par
\textbf{给出全阶论证：}\cref{prop:pdf-support-hadron}（\tagA）；\cref{prop:pdf-timelike-dglap}（\tagA；碎裂函数的乘法可重整性是独立输入，\cref{asm:appN-multren}(Z1)）。\par
\textbf{没有证明：}$\MSbar$ PDF 在某一非微扰标度以上的正定性（\cref{conj:pdf-positivity}，\tagC；在所有标度上非负是错的）；数密度诠释（启发式，\cref{rem:pdf-number-density}）；非定域算符的 BRST 分类（第~24~章讨论其矩约化）。
\end{tcolorbox}

% -----------------------------------------------------------------------------------------------------
\section{定义：夸克、反夸克与胶子分布}\label{sec:pdf-definitions}

部分子模型的图像是：高速强子中部分子间的相互作用因时间膨胀而变慢，硬碰撞“看到”的是一组近似自由、各带加号动量份额 $\xi$ 的部分子。把“份额为 $\xi$ 的夸克数密度”变成定义，Collins 与 Soper 的做法\cite{CollinsSoper1982a}（较早的步骤见\cite{Soper:1976jc,Soper:1979fq}；系统叙述见\cite[第~6.3~节、第~7.5--7.6~节]{Collins2011}）是：用好分量 $\gamma^+\psi$ 计数，用沿光锥的 Wilson 线保证规范不变，再在 $\MSbar$ 中重整化计数中的紫外发散。

\paragraph{约定。}
光锥坐标与 $n_A=(1^+,0^-,\bm 0_{\mathrm T})$、$n_B=(0^+,1^-,\bm 0_{\mathrm T})$ 如\cref{def:prelim-lightfront}：$n_B\cdot v=v^+$，$\sla n_B=\gamma^+$。强子 $H$ 沿 $+z$ 运动，$P=(P^+,\frac{M^2}{2P^+},\bm 0_{\mathrm T})$；分离矢量 $w=(0^+,w^-,\bm 0_{\mathrm T})=w^-n_B$。$D_\mu=\partial_\mu-\ii gA^a_\mu t^a$，Wilson 线 $W[b,a]=P\exp\big(\ii g\int_a^b\dd z\cdot A\big)$，离 $b$ 近的点在左（\cref{def:prelim-wilson}）；沿 $n_B$ 有 $\dd z\cdot A=\dd z^-A^+$，故
\begin{equation}\label{eq:pdf-W}
W[w,0]=P\exp\Big(\ii g\!\int_0^{w^-}\!\dd z^-A^{+}_a(z^-n_B)t^a\Big),\qquad \partial_{w^-}W[w,0]=\ii g\,A^+(w)\,W[w,0].
\end{equation}
态归一为 $\langle P'|P\rangle=2P^+(2\pi)^3\delta(P'^+-P^+)\delta^{(2)}(\bm P'_{\mathrm T}-\bm P_{\mathrm T})$，于是 $\langle P|T^{\mu\nu}|P\rangle=2P^\mu P^\nu$，$\bar u\gamma^\mu u=2p^\mu$。矩阵元一律取连通部分 $\langle\cdots\rangle_{\mathrm c}$，对算符的色指标求和，对靶的色与自旋（极化）求平均。

\begin{definition}[夸克分布]\label{def:pdf-quark}
强子 $H$ 中味道为 $q$ 的夸克分布为
\begin{equation}\label{eq:pdf-fq}
f_{q/H}(\xi)=\int\frac{\dd w^-}{4\pi}\,\ee^{-\ii\xi P^+w^-}\,\big\langle P\big|\bar\psi_q(w)\gamma^+W[w,0]\,\psi_q(0)\big|P\big\rangle_{\mathrm c},
\end{equation}
% VF-PDF-1
$W[w,0]$ 如\cref{eq:pdf-W}。在维数正规化 $D=4-2\varepsilon$ 中\cref{eq:pdf-fq}定义的是裸分布 $f^{\mathrm{bare}}$。$\MSbar$ 分布 $f_{q/H}(\xi;\mu)$ 按如下约定定义：把 $x>0$ 处的各密度 $f_i(x)$（$i$ 遍历各味道的夸克、反夸克与胶子，后两者的定义紧接在后）排成列矢量，
\begin{equation}\label{eq:pdf-Z}
\begin{gathered}
f^{\mathrm{ren}}_i=\sum_jZ_{ij}\otimes f^{\mathrm{bare}}_j,\qquad (Z\otimes f)(x)=\int_x^1\frac{\dd z}{z}\,Z(z)\,f\Big(\frac xz\Big),\\
Z_{qq}(x)=\delta(1-x)-\frac{\as}{2\pi}\frac1\varepsilon P_{qq}(x)+\order{\as^2},
\end{gathered}
\end{equation}
% VF-PDF-4
其中 $Z_{ij}$ 是支撑在 $[0,1]$ 上、用 $\as(\mu)$（$g_0^2=4\pi\as\mu^{2\varepsilon}S_\varepsilon^{-1}Z_g^2$，$S_\varepsilon=(4\pi)^\varepsilon\ee^{-\varepsilon\gamma_E}$，\cref{def:prelim-msbar}）表示时只含 $\varepsilon$ 极点的分布。
\uses{def:prelim-wilson,def:prelim-lightfront,def:prelim-msbar}
\end{definition}

相反的约定 $f^{\mathrm{bare}}=Z'\otimes f^{\mathrm{ren}}$ 有 $Z'_{qq}=\delta+\frac{\as}{2\pi}\frac1\varepsilon P_{qq}+\cdots$，本书只用\cref{eq:pdf-Z}。与靶无关的卷积核 $Z$ 的存在性不是定义的一部分，而是\cref{thm:pdf-moments-local}的内容（证明见\cref{app:appL}）；在那之前我们只对部分子靶使用它。因子 $1/(4\pi)$ 使自由夸克给出 $\delta(1-\xi)$（\cref{prop:pdf-invariance}）。在两个场之间插入末态完备集，用平移 $\langle X|\psi(w)\cdots\rangle\propto\ee^{-\ii(P-p_X)\cdot w}$ 并完成 $w^-$ 积分，得割图表示（\cref{fig:pdf-operator}；图中的 $n$ 即本书的 $n_B$）
\begin{equation}\label{eq:pdf-cut}
f_{q/H}(\xi)=\sum_X\tfrac12\,\delta\big(\xi P^+-(P-p_X)^+\big)\,
\langle P|\bar\psi(0)W[0,\infty n_B]|X\rangle\,\gamma^+\,\langle X|W[\infty n_B,0]\psi(0)|P\rangle ,
\end{equation}
其中用了合成律 $W[w,0]=W[w,\infty n_B]W[\infty n_B,0]$；同样 $W[w,0]=W[w,-\infty n_B]W[-\infty n_B,0]$，故对共线 PDF，Wilson 线伸向未来还是过去无关紧要（对横动量依赖分布则不然，\cref{rem:pdf-gauge-residual}）。由\cref{eq:pdf-cut}，Wilson 线上无胶子附着时算符顶点的 Feynman 规则是 $\frac12\gamma^+\delta(k^+-\xi P^+)$（配合 $\int\dd^4k/(2\pi)^4$）。

%%% BEGIN fig:pdf-operator
\begin{figure}[htbp]
\centering
\tikzset{t1 hadron/.style={draw,solid,line width=1.5pt,line cap=butt}}
\begin{tikzpicture}[fd]
  \coordinate (ol) at (-1.9,2.4); \coordinate (or) at (1.9,2.4);
  \coordinate (bl) at (-1.9,0);   \coordinate (br) at (1.9,0);
  % remnant X crossing the cut (drawn first; the blobs cover the ends)
  \foreach \y in {0.24,0,-0.24} \draw[plain] (bl) ++(0,\y) -- ++(3.8,0);
  \draw[quark] (-1.9,0.43) -- (ol);
  \draw[wilson=0.25] (ol) -- (or);
  \draw[quark] (or) -- (1.9,0.43);
  \draw[t1 hadron] (-3.65,0) -- (bl);
  \draw[t1 hadron] (br) -- (3.65,0);
  \draw[hadron blob] (bl) ellipse[x radius=0.82,y radius=0.45];
  \draw[hadron blob] (br) ellipse[x radius=0.82,y radius=0.45];
  \node[opvertex] at (ol) {}; \node[opvertex] at (or) {};
  \draw[cut] (0,-0.85) -- (0,3.0);
  % labels
  \path[mom={$k$},mom pos=0.62] (-1.9,0.45) -- (ol);
  \path[mom={$k$},mom pos=0.38] (or) -- (1.9,0.45);
  \node[left,font=\small] at (-2.02,0.92) {$k^+=\xi P^+$};
  \node[above left,inner sep=3pt] at (ol) {$\psi(0)$};
  \node[above right,inner sep=3pt] at (or) {$\bar\psi(w)$};
  \node[below,font=\footnotesize,inner sep=4pt] at (-0.95,2.4) {$W[\infty n,0]$};
  \node[below,font=\footnotesize,inner sep=4pt] at (0.95,2.4) {$W[w,\infty n]$};
  \path[mom={$P$},mom sep=0.3cm] (-3.65,0) -- (-2.72,0);
  \path[mom={$P$},mom sep=0.3cm] (2.72,0) -- (3.65,0);
  \node[font=\small] at (0.42,-0.62) {$X$};
\end{tikzpicture}
\caption{夸克分布函数的割图表示。在定义
$f_{q/H}(\xi)=\int\frac{\dd w^-}{4\pi}\,\ee^{-\ii\xi P^+w^-}\langle P|\bar\psi(w)\gamma^+W[w,0]\psi(0)|P\rangle$
（$w=w^-n$，$n$ 为满足 $n\cdot v=v^+$ 的光锥矢量）中插入末态完备集并完成 $w^-$ 积分，得
$f_{q/H}(\xi)=\sum_X\tfrac12\,\delta\big(\xi P^+-(P-p_X)^+\big)\,\langle P|\bar\psi(0)W[0,\infty n]|X\rangle\,\gamma^+\,\langle X|W[\infty n,0]\psi(0)|P\rangle$。
割线左侧为振幅、右侧为共轭振幅，$X$ 为穿过割线的剩余部分。两个 $\otimes$ 为非定域算符的两端，算符顶点的 Feynman 规则为
$\tfrac12\gamma^+\delta(k^+-\xi P^+)$（配合 $\int\dd^4k/(2\pi)^4$，Wilson 线上无胶子附着时）；
双线为 Wilson 线 $W[w,0]=W[w,\infty n]\,W[\infty n,0]$，被末态割线分为振幅中的 $W[\infty n,0]$ 与共轭振幅中的 $W[w,\infty n]$。
箭头表示费米子数（色）流向。}
\label{fig:pdf-operator}
\end{figure}
%%% END fig:pdf-operator

\begin{definition}[反夸克分布]\label{def:pdf-antiquark}
反夸克分布定义为 $f_{\bar q/H}(\xi)\equiv-f_{q/H}(-\xi)$；$x>0$ 时记 $q(x)=f_{q/H}(x)$，$\bar q(x)=f_{\bar q/H}(x)=-f_{q/H}(-x)$。它与用电荷共轭场 $\psi^c=C\bar\psi^{\tp}$（$C=\ii\gamma^2\gamma^0$）及其共轭表示中的 Wilson 线 $W^*$ 直接写出的定义 $\int\frac{\dd w^-}{4\pi}\ee^{-\ii\xi P^+w^-}\langle P|\bar\psi^c(w)\gamma^+W^*[w,0]\psi^c(0)|P\rangle_{\mathrm c}$ 相同。
% VF-PDF-1
\uses{def:pdf-quark}
\end{definition}

等价性：由 $C^{-1}\gamma^\mu C=-\gamma^{\mu\tp}$、$\bar\psi^c=-\psi^{\tp}C^{-1}$，并把 Grassmann 双线性式转置（产生一个负号与一个反对易子），用 $W^{*\tp}[w,0]=W[0,w]$，
\begin{equation}\label{eq:pdf-cconj}
\bar\psi^c(w)\gamma^+W^*[w,0]\psi^c(0)=-\bar\psi(0)\gamma^+W[0,w]\psi(w)+(\text{c 数}).
\end{equation}
c 数与态无关，不出现在连通矩阵元中；再由平移不变性与代换 $w^-\to-w^-$ 即得 $-f_{q/H}(-\xi)$。于是 $f_q$ 是 $[-1,1]$ 上的对象：$\xi>0$ 部分是夸克密度，$\xi<0$ 部分是反号的反夸克密度。

\begin{definition}[胶子分布]\label{def:pdf-gluon}
强子 $H$ 中的胶子分布为
\begin{equation}\label{eq:pdf-fg}
\begin{gathered}
f_{g/H}(\xi)=\sum_{j}\int\frac{\dd w^-}{2\pi\,\xi P^+}\,\ee^{-\ii\xi P^+w^-}\,\big\langle P\big|G^{+j}_a(w)\,W^{(A)}_{ab}[w,0]\,G^{+j}_b(0)\big|P\big\rangle_{\mathrm c},\\
W^{(A)}_{ab}[w,0]=2\Tr\big[t^aW[w,0]\,t^bW[0,w]\big],
\end{gathered}
\end{equation}
% VF-PDF-2
$j$ 遍历 $D-2$ 个横向指标。$f_g$ 是奇分布，$f_g(-\xi)=-f_g(\xi)$（$\xi f_g$ 是偶分布；这也固定了 $1/\xi$ 在 $\xi=0$ 处的含义），$x>0$ 时 $g(x)=f_g(x)$。部分子胶子靶按 $\MSbar$ 约定对 $D-2=2(1-\varepsilon)$ 个极化求平均（色平均 $\frac1{N_c^2-1}$）；改为对 2 个极化求平均会使 DIS 的 $C_2^g\to C_2^g+P_{qg}$、Drell--Yan 的 $\Delta_{qg}\to\Delta_{qg}+P_{qg}$。$f_g$ 与夸克分布一起按\cref{eq:pdf-Z}重整化。
% VF-CONV-6
\uses{def:prelim-wilson,def:pdf-quark}
\end{definition}

在 $A^+=0$ 规范中 $G^{+j}=\partial^+A^j$，因子 $1/(\xi P^+)=1/k^+$ 把算符 $\propto(k^+)^2A^jA^j$ 变成“胶子数”（与标量场的数密度 $\propto k^+|\phi|^2$ 同形）；一般规范中的组合 $k^+A^j-k^jA^+$ 即 Grammer--Yennie 分解中领头的 G 部分（比较\cref{prop:ir-super-leading}与第~9~章的\cref{prop:method-gy-suppression}）。

% -----------------------------------------------------------------------------------------------------
\section{规范不变性、自由场值与支撑}\label{sec:pdf-properties}

\begin{proposition}[规范不变性、推动不变性与自由场值 \tagR]\label{prop:pdf-invariance}
（$R_{\mathrm{def}}$：由算符定义直接推出。）
(a) $f_{q/H}$、$f_{\bar q/H}$、$f_{g/H}$ 在规范变换 $\psi\to U\psi$，$A_\mu\to UA_\mu U^\dagger-\frac{\ii}{g}(\partial_\mu U)U^\dagger$ 下不变。
(b) 它们只依赖于 $\xi$（与 $\mu$），在沿 $z$ 轴的推动 $v^\pm\to\ee^{\pm\eta}v^\pm$ 下不变。
(c) 对动量 $p=(p^+,0^-,\bm 0_{\mathrm T})$ 的无质量自由靶，在 $g^0$ 阶
\begin{equation}\label{eq:pdf-free}
f^{(0)}_{q/q}=\delta(1-\xi),\qquad f^{(0)}_{q/\bar q}=-\delta(1+\xi)\ \ \big(\text{故 }f^{(0)}_{\bar q/\bar q}=\delta(1-\xi)\big),\qquad
f^{(0)}_{g/g}=\delta(1-\xi)-\delta(1+\xi).
\end{equation}
% VF-PDF-1, VF-PDF-2
(d) 胶子的归一由能动张量固定：$T^{++}_g=\sum_jG^{+j}_aG^{+j}_a$ 满足 $\langle k|T^{++}_g|k\rangle=2(k^+)^2$，与 $\langle P|T^{\mu\nu}|P\rangle=2P^\mu P^\nu$ 一致。
非微扰地位：强子态需\cref{asm:intro-wightman}、\cref{asm:intro-composite}（MA0）。
\uses{def:pdf-quark,def:pdf-antiquark,def:pdf-gluon,prop:prelim-wilson-covariance,lem:prelim-boost,asm:intro-wightman,asm:intro-composite}
\end{proposition}

\begin{proof}
(a) 由\cref{prop:prelim-wilson-covariance}，$W[w,0]\to U(w)W[w,0]U^\dagger(0)$，故 $\bar\psi(w)\gamma^+W[w,0]\psi(0)$ 不变。胶子：记 $G=G_at^a$，$G_a(w)W^{(A)}_{ab}[w,0]G_b(0)=2\Tr[G(w)W[w,0]G(0)W[0,w]]$，而 $G\to UGU^\dagger$，迹中的 $U$ 两两相消。
(b) 由\cref{lem:prelim-boost}，好分量 $P_+\psi$（$P_+=\frac12\gamma^-\gamma^+$）按 $\ee^{\eta/2}$ 标度；$\gamma^0\gamma^+=\sqrt2P_+$ 给出 $\bar\psi\gamma^+\psi=\sqrt2\,\psi_+^\dagger\psi_+$，故双线性式按 $\ee^{\eta}$ 标度，而 $w^-\to\ee^{-\eta}w^-$。于是 $M_P(w^-)\equiv\langle P|\bar\psi(w)\gamma^+W[w,0]\psi(0)|P\rangle_{\mathrm c}$ 满足 $M_{\Lambda P}(w^-)=\ee^{\eta}M_P(\ee^{\eta}w^-)$（协变归一在推动下不变），代入\cref{eq:pdf-fq}并换元 $w'=\ee^\eta w^-$ 即得 $f$ 不变。胶子同理（$G^{+j}G^{+j}$ 按 $\ee^{2\eta}$ 标度，$1/(\xi P^+)$ 补偿一个 $\ee^{\eta}$）。
(c) 在 $g^0$ 阶 $W=1$，$\psi=\int[b\,u\,\ee^{-\ii p\cdot x}+d^\dagger v\,\ee^{\ii p\cdot x}]$。夸克靶的连通矩阵元为 $\bar u\gamma^+u\,\ee^{\ii p^+w^-}=2p^+\ee^{\ii p^+w^-}$（对每个螺旋度 $\bar u\gamma^+u=2p^+$，显式验证见下例），故 $f^{(0)}_{q/q}=\frac{2p^+}{4\pi}\,2\pi\delta((1-\xi)p^+)=\delta(1-\xi)$。反夸克靶：$\bar\psi(w)$ 中的 $d\,\bar v\,\ee^{-\ii p\cdot w}$ 与 $\psi(0)$ 中的 $d^\dagger v$ 缩并，化为正规序时 $dd^\dagger\to-d^\dagger d$，矩阵元为 $-\bar v\gamma^+v\,\ee^{-\ii p^+w^-}=-2p^+\ee^{-\ii p^+w^-}$，故 $f^{(0)}_{q/\bar q}=-\delta(1+\xi)$。胶子靶（$\bm k_{\mathrm T}=0$）：$\langle0|G^{+j}_a(x)|k,\lambda,b\rangle=-\ii k^+\epsilon^j_\lambda\delta_{ab}\ee^{-\ii k\cdot x}$；连通矩阵元含“$0$ 处湮灭、$w$ 处产生”及其相反两项，对极化平均后 $\sum_j|\epsilon^j|^2\to1$，得 $(k^+)^2(\ee^{\ii k^+w^-}+\ee^{-\ii k^+w^-})$，于是 $f^{(0)}_{g/g}=\frac{(k^+)^2}{2\pi\xi k^+}\frac{2\pi}{k^+}[\delta(\xi-1)+\delta(\xi+1)]=\delta(1-\xi)-\delta(1+\xi)$（$\xi=-1$ 处 $1/\xi=-1$）。
(d) 同一胶子态中 $\sum_jG^{+j}G^{+j}$ 的矩阵元是上式在 $w^-=0$ 处的值 $2(k^+)^2$；另一方面 $T^{++}_g=-G^{+\alpha}_aG^{+}{}_{\alpha a}$ 是规范不变能动张量 $-G^{\mu\alpha}_aG^\nu{}_{\alpha a}+\frac14g^{\mu\nu}G^2$ 的 $++$ 分量（$G^{++}=0$，$G^+{}_-=0$，只剩横向求和），$\langle k|T^{\mu\nu}|k\rangle=2k^\mu k^\nu$ 给出同一值。
\end{proof}

\begin{example}[自由场值]\label{ex:pdf-free-field}
Dirac 表示中，$p$ 沿 $+z$、$E=p^+/\sqrt2$，螺旋度 $h$ 的旋量 $u=\sqrt E(\chi_h,2h\chi_h)^{\tp}$；由 $\gamma^0\gamma^+=(1+\gamma^0\gamma^3)/\sqrt2$，$\bar u\gamma^+u=\frac{E}{\sqrt2}(2+2(2h)^2)=2p^+$，与 $h$ 无关。割图表示\cref{eq:pdf-cut}中夸克靶只有 $X=$ 真空一项，给出 $\frac12\delta(\xi p^+-p^+)\cdot2p^+=\delta(1-\xi)$；反夸克靶的负号来自费米统计而非旋量（用\cref{def:pdf-antiquark}中电荷共轭场的定义直接得 $f^{(0)}_{\bar q/\bar q}=\delta(1-\xi)$）。各自由分布的矩 $1$、$(-1)^n$、$1+(-1)^n$ 与定域算符 $\bar\psi\gamma^+(\ii\partial^+)^{n-1}\psi$、$\sum_jG^{+j}(\ii\partial^+)^{n-2}G^{+j}$ 在同样态中的值除以 $2(p^+)^n$ 逐一相等（\cref{thm:pdf-moments-local}的自由场特例）；自由胶子的 $\int_0^1x\,g=1$ 即动量求和规则（\cref{prop:pdf-sum-rules}）。
\uses{prop:pdf-invariance}
\end{example}

\begin{theorem}[支撑（微扰论）\tagR]\label{thm:pdf-support-pt}
（$R_{\mathrm{PT}}$：在重整化微扰论中对红外正规化的部分子态严格。）
设靶为红外正规化的在壳部分子 $j$。在微扰论每一阶、对每个割图，裸分布 $f_{q/j}(\xi)$ 与 $f_{g/j}(\xi)$ 在 $|\xi|>1$ 时为零，即它们是支撑在 $[-1,1]$ 上的分布（等价地，$x>0$ 处的各密度支撑在 $[0,1]$ 上）。$\MSbar$ 重整化保持这一性质，因为逐阶由极点部分构造的抵消项 $Z_{ij}(z)$ 支撑在 $0\le z\le1$ 上（$Z(\xi/\zeta)$ 在 $\xi>\zeta$ 时为零）。（本定理只涉及部分子态。）
\uses{def:pdf-quark,def:pdf-gluon,def:ir-landau-setup}
\end{theorem}

\begin{proofsketch}
(1) 微扰论中\cref{eq:pdf-fq}的矩阵元是普通（非编时）乘积；在 $[\bar\psi W](w)$ 与 $[W\psi](0)$ 之间插入末态完备集得割图求和\cref{eq:pdf-cut}。每个割图中穿过割线的都是在壳正能量粒子，各自 $p^+=(E+p^3)/\sqrt2\ge0$，故 $p_X^+\ge0$，$\delta$ 函数给出 $\xi p^+=p^+-p_X^+\le p^+$。Wilson 线只是源，不携带在壳末态。(2) 由\cref{eq:pdf-cconj}，$f_q(\xi)=-f_{\bar q}(-\xi)$，对电荷共轭场的分布重复 (1) 得 $\xi\ge-1$；反对易子的 c 数与 $\xi$ 无关，不进入连通矩阵元。(3) 胶子以伴随 Wilson 线同样处理。(4) 对部分子靶，$Z^{(k)}=-\mathrm{PP}\big[f^{(k)}_{\mathrm{bare}}+\sum_{l<k}Z^{(l)}\otimes f^{(k-l)}_{\mathrm{bare}}\big]$（PP 为 $\varepsilon$ 的极点部分）；支撑在 $[0,1]$ 上的分布的 Mellin 卷积仍支撑在 $[0,1]$ 上，取 Laurent 系数不扩大支撑，归纳即得。这是谱条件论证在图上的版本；用光锥序微扰论的表述见\cite[第~6.9.3~节，第~191~页]{Collins2011}（该处指出论证是非微扰的），对多部分子关联函数的推广见\cite{Jaffe1983}（据该文摘要，每个量子的动量分数都在 $0$ 与 $1$ 之间），时间排序与支撑的关系见\cite{Diehl:1998sm}（据该文摘要）。完整论证见上述文献。
\end{proofsketch}

\begin{proposition}[支撑（强子态）\tagA]\label{prop:pdf-support-hadron}
对强子态，$f_{q/H}(\xi)=0$ 与 $f_{g/H}(\xi)=0$ 对 $|\xi|>1$ 成立：(a) 在 $\bar\psi(w)W$ 与 $W\psi(0)$ 之间插入完备态，$w^-$ 积分把中间态的加号动量固定为 $p_X^+=(1-\xi)P^+$，谱条件给出 $p_X^+\ge0$，故 $\xi\le1$；(b) 反夸克部分由光类间隔上的反对易子给出（扣除与态无关的发散常数），由此 $\xi\ge-1$。
非微扰地位：\cref{asm:intro-wightman}、\cref{asm:intro-composite}（MA0）。
\uses{def:pdf-quark,def:pdf-antiquark,thm:pdf-support-pt,asm:intro-wightman,asm:intro-composite}
\end{proposition}

\begin{proofsketch}
(1) 由\cref{asm:intro-wightman}(W2)，联合谱在闭前向锥内，故任何中间态 $p_X^0\ge|\bm p_X|$，$p_X^+\ge0$。(2) 插入 $\sum_X|X\rangle\langle X|$，与\cref{thm:pdf-support-pt}第 (1) 步相同地得 $\xi\le1$。(3) 写 $\bar\psi(w)\Gamma\psi(0)=-\psi(0)^{\tp}\Gamma^{\tp}\bar\psi^{\tp}(w)+\{\cdot,\cdot\}$；反对易子在光类间隔上是与态无关的发散常数，从连通矩阵元中扣除；对剩下的项插入完备态得 $\xi\ge-1$。(4) 胶子同理。论证是非微扰的\cite[第~6.9.3~节，第~191~页]{Collins2011}，微扰论中的对应是\cref{thm:pdf-support-pt}。
\end{proofsketch}

\paragraph{地位与缺口}
谱论证把\emph{重整化的}光线算符当作可以在两半之间插入完备态的算符乘积；而重整化作用在整个双局域算符上，并不是两个“半算符”各自重整化之积。“对态求和”与“重整化”的交换是未经证明的一步，它使本命题停留在 \tagA；所依赖的假设是\cref{asm:intro-wightman}(W2) 与\cref{asm:intro-composite}(C3)（光线算符在强子态之间作为缓增分布存在），这一步隐含在\cite[第~6.9.3~节]{Collins2011}的谱论证中。此缺口与幂次计数无关，不涉及 E-S 至 E-W 型的误差估计。$\xi=0$ 一点需要小心：扭度 3 的分布可以含 $\delta(x)$ 项，目前不知道扭度 2 的非极化 PDF 有这类项。若能非微扰地构造重整化光线算符并证明其谱表示，结论即可在\cref{asm:intro-wightman}、\cref{asm:intro-composite}之下升级为 $[R,R_{\mathrm{cond}}]$。第二种论证见\cref{rem:pdf-support-jaffe}。

\begin{remark}[支撑的第二种论证：$k^-$ 极点]\label{rem:pdf-support-jaffe}
写 $f(\xi)=\int\frac{\dd^4k}{(2\pi)^4}\delta(k^+-\xi p^+)\TrD[\frac12\gamma^+\mathcal A(k,p)]$，$\mathcal A$ 是部分子—靶前向振幅。携带加号动量 $\ell^+$ 的传播子在 $\ell^-=(\bm\ell_{\mathrm T}^2+m^2-\ii0)/(2\ell^+)$ 处有极点，$\ell^+>0$ 时在实轴下方，$\ell^+<0$ 时在上方。当 $\xi>1$ 时，靶取走一个部分子后的中间态总加号动量 $(1-\xi)p^+$ 为负（$\xi<-1$ 时交换夸克与反夸克的角色），完成其余回路的减号分量积分后，$\mathcal A$ 作为 $k^-$ 的函数的奇点全在实轴同一侧；被积函数在 $|k^-|\to\infty$ 时足够快地衰减，向另一侧闭合围道即得零\cite{Jaffe1983}（据该文摘要，论证要求中间强子态在壳且没有相对于靶或中间态半不连通的图）。微扰论中它与\cref{thm:pdf-support-pt}等价，两者都表述“中间态加号动量非负”。
\uses{thm:pdf-support-pt}
\end{remark}

% -----------------------------------------------------------------------------------------------------
\section{数密度诠释及其局限（启发式）；光锥规范的剩余规范自由}\label{sec:pdf-interpretation}

算符定义的直观读法来自光锥规范 $A^+=0$（此时 $W[w,0]=1$）；它只是启发式的，并且正是“$\MSbar$ PDF 是概率、故为正”这一误解的来源。

\begin{remark}[数密度诠释（启发式）]\label{rem:pdf-number-density}
在 $A^+=0$ 规范中按光前产生、湮灭算符（$\{b,b^\dagger\}=(2\pi)^32k^+\delta^{(3)}$）展开好分量，用 $\bar\psi\gamma^+\psi=\sqrt2\psi_+^\dagger\psi_+$、$u_+^\dagger u_+=\sqrt2k^+$，对 $0<\xi<1$ 得
\begin{equation}\label{eq:pdf-number}
f_{q/H}(\xi)=\sum_{s,c}\int\frac{\dd^{2}\bm k_{\mathrm T}}{2\xi\,(2\pi)^3}\,
\frac{\big\|b_{s,c}(\xi P^+,\bm k_{\mathrm T})|P\rangle\big\|^2}{\langle P|P\rangle}\ \ge0
\end{equation}
（对自由夸克，比值为 $(2\pi)^32p^+\delta^{(3)}(k-p)$，给出 $\delta(1-\xi)$），即“份额为 $\xi$ 的夸克数密度”\cite[第~6.7~节、第~6.9.1~节]{Collins2011}。然而：(i) 相互作用理论中 $\bm k_{\mathrm T}$ 积分在紫外对数发散，正是 $\MSbar$ 所要去掉的发散，重整化后正性不复存在——Collins 写道：「“this particular result will not hold exactly in renormalizable theories, because of the need for renormalization”」\cite[第~6.9.1~节]{Collins2011}，即因为需要重整化，该结果在可重整理论中不严格成立；(ii) $A^+=0$ 规范的伪 $1/k^+$ 奇异性需要处方（\cref{rem:pdf-gauge-residual}）；(iii) 光前量子化有 $k^+=0$ 的零模问题；(iv) 光锥无穷远处的横向规范连接在 $A^+=0$ 规范中并不消失\cite{JiYuan2002,BelitskyJiYuan2003}，对共线 PDF 它们因两端横向位置相同而相消，对横动量依赖分布则不然。所以\cref{eq:pdf-number}是关于\emph{裸的、紫外截断的}量的启发式陈述，不是关于 $\MSbar$ PDF 的陈述；后者的支撑与求和规则另有证明（\cref{thm:pdf-support-pt}、\cref{prop:pdf-sum-rules}），正定性是\cref{conj:pdf-positivity}。
\end{remark}

\begin{remark}[光锥规范的剩余规范自由与 $1/k^+$ 处方]\label{rem:pdf-gauge-residual}
$A^+=0$ 并未完全固定规范：保持它的变换与 $x^-$ 无关，作用在 $x^-=\pm\infty$ 处；相应地，传播子在 $k^+=0$ 的奇异性需要处方（主值、$k^+\pm\ii0$、Mandelstam--Leibbrandt 等），每种处方等价于对 $\bm A_{\mathrm T}$ 在 $x^-=\pm\infty$ 的一种边界条件。共线 PDF 的 Wilson 线是横向位置相同的\emph{有限}类光线段，$W[w,0]=W[w,\pm\infty n_B]W[\pm\infty n_B,0]$，无穷远横向连接平凡，故结果与处方无关；横动量依赖分布的订书钉形 Wilson 线两端横向位置不同，结果依赖于边界条件（第~14~章）。常见论证“用规范变换把 $W$ 变成 $1$”中的 $U(x)=W[\infty n_B,x]$ 是依赖于场的算符而非整体规范变换；结论仍成立，只因它用到的仅是合成律。
\end{remark}

% -----------------------------------------------------------------------------------------------------
\section{矩与定域算符}\label{sec:pdf-moments}

第~7~章的算符乘积展开给出结构函数的\emph{矩}：每个偶数矩等于 Wilson 系数乘以扭度 2 定域算符的约化矩阵元 $A^{(n)}_i$（\cref{thm:dis-moments}）。与本章定义的桥梁是“PDF 的 $n$ 阶矩 $=A^{(n)}$”：对裸算符，它是光线算符在 $w^-=0$ 处的 Taylor 展开；对重整化的量，它要求 PDF 重整化核的矩恰为定域算符的重整化常数。

\begin{definition}[光线算符]\label{def:pdf-lightray}
称 $O_q(w^-)\equiv\bar\psi_q(w)\gamma^+W[w,0]\psi_q(0)$（$w=w^-n_B$）为夸克光线算符（light-ray operator），$O_g(w^-)\equiv\sum_jG^{+j}_a(w)W^{(A)}_{ab}[w,0]G^{+j}_b(0)$ 为胶子光线算符。其前向矩阵元与分布互为 Fourier 变换：
\begin{equation}\label{eq:pdf-ioffe}
\begin{gathered}
M_q(w^-)\equiv\langle P|O_q(w^-)|P\rangle_{\mathrm c}=2P^+\!\int\!\dd\xi\,\ee^{\ii\xi P^+w^-}f_q(\xi),\\
M_g(w^-)\equiv\langle P|O_g(w^-)|P\rangle_{\mathrm c}=(P^+)^2\!\int\!\dd\xi\,\ee^{\ii\xi P^+w^-}\xi f_g(\xi).
\end{gathered}
\end{equation}
由推动不变性，$M_q/P^+$ 与 $M_g/(P^+)^2$ 只通过 $P\cdot w=P^+w^-$ 依赖于 $w^-$。
\uses{def:pdf-quark,def:pdf-gluon}
\end{definition}

\cref{eq:pdf-ioffe}是\cref{eq:pdf-fq}、\cref{eq:pdf-fg}的逆变换。$P\cdot w$ 在光锥上的作用相当于 Ioffe 时间；第~22~章的准分布与赝分布使用类空分离的类比（\cref{def:lattice-quasi}）。

\begin{lemma}[Wilson 线顶点的矩 = 定域算符顶点 \tagR]\label{lem:pdf-wilson-moments}
（$R_{\mathrm{def}}$：由算符定义直接推出。）
设算符顶点不注入动量（流出夸克的动量等于流入夸克与全部附着胶子的动量之和）。对 $n\ge1$、$m\ge0$，$\int\frac{\dd w^-}{4\pi}\ee^{-\ii\xi P^+w^-}O_q(w^-)$ 的 $m$ 胶子顶点的 $n$ 阶矩 $\int\dd\xi\,\xi^{n-1}(\cdots)$ 等于 $(2(P^+)^n)^{-1}$ 乘以 $\bar\psi\gamma^+(\ii D^+)^{n-1}\psi$ 的 $m$ 胶子顶点（$m\ge n$ 时两者为零）。单胶子情形（胶子动量 $q$、色 $a$、指标 $\mu$ 流入；夸克 $k_1$ 流入、$k_2=k_1+q$ 流出）：
\begin{equation}\label{eq:pdf-vertex1}
\frac{g\,t^an_B^\mu\gamma^+}{2q^+}\Big[\delta(\xi P^+-k_2^+)-\delta(\xi P^+-k_1^+)\Big]\ \xrightarrow{\ n\ \text{阶矩}\ }\ \frac{g\,t^an_B^\mu\gamma^+}{2(P^+)^n}\sum_{j=0}^{n-2}(k_2^+)^j(k_1^+)^{n-2-j}.
\end{equation}
% VF-PDF-12
一般地，对固定的附着次序（色因子 $t^{a_m}\cdots t^{a_1}$），两边都等于 $g^mn_B^{\mu_m}\cdots n_B^{\mu_1}t^{a_m}\cdots t^{a_1}\gamma^+$ 乘以均差 $x^{n-1}[\kappa_0,\dots,\kappa_m]=h_{n-1-m}(\kappa_0,\dots,\kappa_m)$，$\kappa_i=k_1^++q_1^++\dots+q_i^+$，$h_r$ 为 $r$ 次完全齐次对称多项式。
\uses{def:pdf-lightray,prop:prelim-wilson-rules,lem:appL-divided-differences}
\end{lemma}

\begin{proof}
$m=0$：$\int\frac{\dd w^-}{4\pi}\ee^{-\ii(\xi P^+-k^+)w^-}=\frac12\delta(\xi P^+-k^+)$，其 $n$ 阶矩 $\frac12(k^+)^{n-1}/(P^+)^n$ 对应零胶子顶点 $\gamma^+(k^+)^{n-1}$。$m=1$：由\cref{eq:pdf-W}与 $A(x)=\int\frac{\dd^4q}{(2\pi)^4}\ee^{-\ii q\cdot x}\tilde A(q)$（$q$ 流入 Wilson 线，\cref{prop:prelim-wilson-rules}），$W[w,0]$ 的单胶子项为 $\ii g\,t^a\tilde A^{a+}\int_0^{w^-}\dd z\,\ee^{-\ii q^+z}=g\,t^a\tilde A^{a+}(1-\ee^{-\ii q^+w^-})/q^+$（$q^+=0$ 处极点可去）；乘上 $\bar\psi(w)$ 的 $\ee^{\ii k_2^+w^-}$，$w^-$ 依赖为 $\frac{g}{q^+}(\ee^{\ii k_2^+w^-}-\ee^{\ii k_1^+w^-})$，Fourier 变换即\cref{eq:pdf-vertex1}左边，其矩为 $\frac{g\,t^an_B^\mu\gamma^+}{2(P^+)^n}\frac{(k_2^+)^{n-1}-(k_1^+)^{n-1}}{k_2^+-k_1^+}$，即右边。另一方面 $\ii D^+=\ii\partial^++gA^+$，$(\ii D^+)^{n-1}$ 中含一个 $gA^+$ 的项为 $\sum_j(\ii\partial^+)^jgA^+(\ii\partial^+)^{n-2-j}$；$\ii\partial^+$ 给出其右侧全部场的加号动量之和，即 $A^+$ 右侧为 $k_1^+$、左侧为 $k_2^+$，定域顶点恰与之相同。一般 $m$ 见\cref{lem:appL-divided-differences}：路径排序积分是 $\ee^{\ii xw^-}$ 的均差（Hermite--Genocchi 公式），取矩化为 $x^{n-1}$ 的均差，而定域顶点按插入位置展开恰为完全齐次对称多项式。
\end{proof}

\begin{theorem}[PDF 的矩 = 扭度 2 定域算符矩阵元 \tagR]\label{thm:pdf-moments-local}
（$R_{\mathrm{PT}}$：在重整化微扰论中对红外正规化的部分子态严格。）
对 $\MSbar$ 重整化的量与 $n\ge1$，
\begin{equation}\label{eq:pdf-moments}
\int_{-1}^1\dd\xi\,\xi^{n-1}f_q(\xi;\mu)=\frac{\langle P|[\bar\psi\gamma^+(\ii D^+)^{n-1}\psi]_\mu|P\rangle}{2(P^+)^n}=A^{(n)}_q(\mu)=\int_0^1\dd x\,x^{n-1}\big[q+(-1)^n\bar q\big](x;\mu),
\end{equation}
\begin{equation}\label{eq:pdf-moments-g}
\int_0^1\dd x\,x^{n-1}g(x;\mu)=\frac{\langle P|[O_g^{(n)}]_\mu|P\rangle}{2(P^+)^n}=A^{(n)}_g(\mu)\quad(n\ \text{偶}),\qquad O_g^{(n)}=\sum_jG^{+j}(\ii D^+)^{n-2}G^{+j},
\end{equation}
% VF-PDF-11
奇数 $n$ 时 $\langle P|O_g^{(n)}|P\rangle=0=\int_{-1}^1\xi^{n-1}f_g\,\dd\xi$。这里 $[\cdot]_\mu$ 是 $\MSbar$ 重整化的扭度 2 定域算符（\cref{def:dis-twist}），$\langle P|O_i^{\mu_1\cdots\mu_n}|P\rangle=2A^{(n)}_i(P^{\mu_1}\cdots P^{\mu_n}-\text{迹})$。同时，\cref{eq:pdf-Z}中与靶无关的卷积核 $Z$ 存在，其 Mellin 矩等于定域算符的重整化矩阵：$\int_0^1\dd z\,z^{N-1}Z(z)=Z_O(N)$（偶数 $N$ 在 $(q_k+\bar q_k,\,g)$ 扇区，奇数 $N$ 在 $q_k-\bar q_k$ 扇区）。
非微扰地位：对强子态需\cref{asm:intro-wightman}、\cref{asm:intro-composite}(C3)（MA0）。
\uses{lem:pdf-wilson-moments,thm:prelim-composite-renorm,def:dis-twist,thm:pdf-support-pt,prop:appL-bare,prop:appL-renormalized,asm:intro-wightman,asm:intro-composite}
\end{theorem}

\begin{proof}
证明见\cref{app:appL}，要点有三。(1) 裸恒等式（\cref{prop:appL-bare}）：由支撑（\cref{thm:pdf-support-pt}），$M_q(w^-)$ 是整函数，矩是 $w^-=0$ 处的导数；由\cref{eq:pdf-W}，$\partial_{w^-}[\bar\psi(w)W[w,0]]=[\bar\psi\overleftarrow{D}{}^+](w)W[w,0]$，迭代给出裸定域算符，前向矩阵元中全导数为零，导数可移到右边——在顶点层次即\cref{lem:pdf-wilson-moments}；胶子用偶分布 $\xi f_g$。(2) 重整化（\cref{prop:appL-renormalized}）：由\cref{thm:prelim-composite-renorm}，红外正规化的部分子态中扭度 2 定域算符在 $\{O_q,O_g\}$ 内乘法重整化；在 $x$ 空间对部分子靶逐阶极小减除得 $Z(x)$，由极小减除的唯一性得 $Z(N)=Z_O(N)$。(3) 对任意态，$Z\otimes f_{\mathrm{bare}}$ 的每个极点系数是支撑在 $[-1,1]$ 上、各阶矩为零的分布，其 Fourier--Laplace 变换是 Taylor 系数全为零的整函数，故为零。
\end{proof}

三点说明。(i) 单圈时 (1)(2) 可精确到 $\varepsilon$ 逐图验证（\cref{ex:pdf-moments-exact}）。(ii) 左作用与右作用的协变导数只在\emph{前向}矩阵元中相等；对广义部分子分布（第~19~章\cref{def:excl-gpd}）两者相差全导数项。(iii) 本定理之后，\cref{thm:dis-moments}中的 $A^{(n)}_i$ 就是本章 PDF 的矩；对强子态，矩关系是\cref{asm:intro-composite}(C3) 的一部分。

% -----------------------------------------------------------------------------------------------------
\section{重整化与 DGLAP 方程}\label{sec:pdf-dglap}

DGLAP 方程常被“推导”为截面与因子化标度无关的推论；若再用它“证明”截面与标度无关，就构成循环论证。本书的次序是：先由算符重整化得到 DGLAP（本节），再证明因子化（第~9~至~12~章），截面的标度无关性只是推论。

%%% BEGIN fig:splitting-vertices
\begin{figure}[htbp]
\centering
\begin{tikzpicture}[fd]
  % ---------- (a) q -> q g ----------
  \begin{scope}[shift={(0,0)}]
    \coordinate (V) at (0,0);
    \draw[quark] (-1.6,0) -- (V);
    \draw[quark] (V) -- (1.45,0.95);
    \draw[gluon] (V) -- (1.45,-0.95);
    \node[vertex] at (V) {};
    \path[mom={$p$}] (-1.6,0) -- (V);
    \path[mom={$zp$}] (V) -- (1.45,0.95);
    \path[mom'={$(1-z)p$}] (V) -- (1.45,-0.95);
    \node[left] at (-1.6,0) {$q$};
    \node[right] at (1.45,0.95) {$q$};
    \node[right] at (1.45,-0.95) {$g$};
    \node[font=\small,align=center,anchor=north] at (0,-1.55)
      {(a)\enspace$q\to qg$\\[2pt]$P_{qq}(z)=C_F\Big[\dfrac{1+z^2}{1-z}\Big]_+$};
  \end{scope}
  % ---------- (b) q -> g q ----------
  \begin{scope}[shift={(7.4,0)}]
    \coordinate (V) at (0,0);
    \draw[quark] (-1.6,0) -- (V);
    \draw[gluon] (V) -- (1.45,0.95);
    \draw[quark] (V) -- (1.45,-0.95);
    \node[vertex] at (V) {};
    \path[mom={$p$}] (-1.6,0) -- (V);
    \path[mom={$zp$}] (V) -- (1.45,0.95);
    \path[mom'={$(1-z)p$}] (V) -- (1.45,-0.95);
    \node[left] at (-1.6,0) {$q$};
    \node[right] at (1.45,0.95) {$g$};
    \node[right] at (1.45,-0.95) {$q$};
    \node[font=\small,align=center,anchor=north] at (0,-1.55)
      {(b)\enspace$q\to gq$\\[2pt]$P_{gq}(z)=C_F\,\dfrac{1+(1-z)^2}{z}$};
  \end{scope}
  % ---------- (c) g -> q qbar (antiquark: fermion arrow into the vertex) ----------
  \begin{scope}[shift={(0,-4.6)}]
    \coordinate (V) at (0,0);
    \draw[gluon] (-1.6,0) -- (V);
    \draw[quark] (V) -- (1.45,0.95);
    \draw[quark] (1.45,-0.95) -- (V);
    \node[vertex] at (V) {};
    \path[mom={$p$}] (-1.6,0) -- (V);
    \path[mom={$zp$}] (V) -- (1.45,0.95);
    \path[mom'={$(1-z)p$}] (V) -- (1.45,-0.95);
    \node[left] at (-1.6,0) {$g$};
    \node[right] at (1.45,0.95) {$q$};
    \node[right] at (1.45,-0.95) {$\bar q$};
    \node[font=\small,align=center,anchor=north] at (0,-1.55)
      {(c)\enspace$g\to q\bar q$\\[2pt]$P_{qg}(z)=T_F\big[z^2+(1-z)^2\big]$};
  \end{scope}
  % ---------- (d) g -> g g ----------
  \begin{scope}[shift={(7.4,-4.6)}]
    \coordinate (V) at (0,0);
    \draw[gluon] (-1.6,0) -- (V);
    \draw[gluon] (V) -- (1.45,0.95);
    \draw[gluon] (V) -- (1.45,-0.95);
    \node[vertex] at (V) {};
    \path[mom={$p$}] (-1.6,0) -- (V);
    \path[mom={$zp$}] (V) -- (1.45,0.95);
    \path[mom'={$(1-z)p$}] (V) -- (1.45,-0.95);
    \node[left] at (-1.6,0) {$g$};
    \node[right] at (1.45,0.95) {$g$};
    \node[right] at (1.45,-0.95) {$g$};
    \node[font=\small,align=center,anchor=north] at (0,-1.55)
      {(d)\enspace$g\to gg$\\[2pt]$\begin{aligned}P_{gg}(z)={}&2C_A\Big[\dfrac{z}{(1-z)_+}+\dfrac{1-z}{z}+z(1-z)\Big]\\[-2pt]&+\dfrac{\beta_0}{2}\,\delta(1-z)\end{aligned}$};
  \end{scope}
\end{tikzpicture}
\caption{领头阶分裂函数对应的四种 $1\to2$ 分裂。动量为 $p$ 的母部分子 $j$（左）分裂为携带加号动量 $zp^+$ 的部分子 $i$（右上）
与携带 $(1-z)p^+$ 的部分子（右下）；动量箭头旁的 $zp$、$(1-z)p$ 只标出沿共线方向的动量分数（略去横动量）。
$P_{ij}(z)$ 是 DGLAP 方程 $\partial f_i/\partial\ln\mu^2=\frac{\as}{2\pi}\sum_jP_{ij}\otimes f_j$ 中 $j\to i$ 的领头阶核；
$[\,\cdot\,]_+$ 为加号分布，$C_F[(1+z^2)/(1-z)]_+=C_F[(1+z^2)/(1-z)_++\tfrac32\delta(1-z)]$，$\beta_0=\tfrac{11}{3}C_A-\tfrac43T_Fn_f$。
(c) 中 $P_{qg}$ 对每一味夸克与反夸克分别计入；反夸克线上的箭头（费米子数流向）指向顶点，其动量 $(1-z)p$ 则向外。}
\label{fig:splitting-vertices}
\end{figure}
%%% END fig:splitting-vertices

\begin{theorem}[由算符重整化得到 DGLAP 方程 \tagR]\label{thm:pdf-dglap}
（$R_{\mathrm{PT}}$：在重整化微扰论中对红外正规化的部分子态严格。）
$\MSbar$ 分布满足
\begin{equation}\label{eq:pdf-dglap}
\frac{\partial f_i(x;\mu)}{\partial\ln\mu^2}=\sum_j\big(\mathcal P_{ij}\otimes f_j\big)(x),\qquad
\mathcal P_{ij}=\frac{\as}{2\pi}P_{ij}+\order{\as^2},
\end{equation}
% VF-PDF-7 (normalization), VF-PDF-4
$P_{ij}$ 为领头阶核（显式形式见紧接的命题）。把\cref{eq:pdf-Z}中的 $Z$ 写成 $Z=\mathbb 1+\sum_{k\ge1}Z^{(k)}/\varepsilon^k$，则 $\mathcal P=-\as\,\partial Z^{(1)}/\partial\as$，单圈 $Z^{(1)}=-\frac{\as}{2\pi}P+\order{\as^2}$；高阶极点 $Z^{(k)}$（$k\ge2$）由 $Z^{(1)}$ 与 $\beta$ 函数完全确定，使 $\mathcal P$ 有限。
非微扰地位：强子态需\cref{asm:intro-composite}（MA0；部分子态无需）。
\uses{def:pdf-quark,def:pdf-gluon,thm:pdf-moments-local,prop:qcd-rge,prop:prelim-beta,asm:intro-composite}
\end{theorem}

\begin{proof}
(1) 由\cref{thm:pdf-moments-local}，$f^{\mathrm{ren}}=Z\otimes f^{\mathrm{bare}}$，$Z$ 是与靶无关的味道矩阵值卷积核，只含极点，且只通过 $\as(\mu)$ 依赖于 $\mu$（极小减除的抵消项不含显式的 $\mu$，\cref{prop:qcd-rge}）；裸算符与 $g_0$ 不含 $\mu$，故 $f^{\mathrm{bare}}$ 与 $\mu$ 无关。(2) $D$ 维中 $\mu^2\frac{\dd\as}{\dd\mu^2}=-\varepsilon\as+\beta(\as)$，$\beta=-\frac{\beta_0}{4\pi}\as^2+\order{\as^3}$（\cref{prop:prelim-beta}）。(3) 于是 $\partial_{\ln\mu^2}f^{\mathrm{ren}}=\frac{\dd Z}{\dd\ln\mu^2}\otimes f^{\mathrm{bare}}=\mathcal P\otimes f^{\mathrm{ren}}$，$\mathcal P\equiv\frac{\dd Z}{\dd\ln\mu^2}\otimes Z^{-1}$，$Z^{-1}$ 为卷积代数中的逆（$Z=\delta\,\mathbb 1+\order{\as}$ 作为形式幂级数可逆），矩阵次序如所写。(4) 展开：
\[
\mathcal P=(-\varepsilon\as+\beta)\frac{\partial Z}{\partial\as}\otimes Z^{-1}
=-\as\frac{\partial Z^{(1)}}{\partial\as}+\frac1\varepsilon\Big[-\as\frac{\partial Z^{(2)}}{\partial\as}+\as\frac{\partial Z^{(1)}}{\partial\as}\otimes Z^{(1)}+\beta\frac{\partial Z^{(1)}}{\partial\as}\Big]+\order{\varepsilon^{-2}}.
\]
(5) 对每个靶，$f^{\mathrm{ren}}$ 与其 $\ln\mu^2$ 导数有限（每阶系数是 $\ln\mu^2$ 的多项式）；取一组部分子靶使 $f^{\mathrm{ren}}_{i/j}=\delta_{ij}\delta(1-x)+\order{\as}$ 可逆，则 $\mathcal P=(\partial_{\ln\mu^2}f^{\mathrm{ren}})\otimes(f^{\mathrm{ren}})^{-1}$ 有限。故 (4) 中 $\varepsilon$ 的负幂系数逐个为零（'t Hooft 一致性条件，如 $\as\partial_{\as}Z^{(2)}=(\as\partial_{\as}Z^{(1)})\otimes Z^{(1)}+\beta\,\partial_{\as}Z^{(1)}$），有限部分给出 $\mathcal P=-\as\partial Z^{(1)}/\partial\as$。(6) 由\cref{eq:pdf-Z}，$Z^{(1)}=-\frac{\as}{2\pi}P+\order{\as^2}$，故 $\mathcal P=\frac{\as}{2\pi}P+\order{\as^2}$；两圈的一致性条件给出 $Z^{(2)}=(\frac{\as}{2\pi})^2[\frac12P\otimes P+\frac{\beta_0}4P]+\order{\as^3}$。(7) 推导只用到算符定义与其重整化，没有用任何因子化定理；对强子态，$f^{\mathrm{ren}}$ 的存在由\cref{asm:intro-composite}(C3) 保证，(3) 的代数不变。
\end{proof}

\cref{eq:pdf-dglap}的物理图像见\cref{fig:dglap-ladder}：在物理规范中，领头对数来自横动量强序的梯形图。这只是图像，定理是\cref{thm:pdf-dglap}。

%%% BEGIN fig:dglap-ladder
\begin{figure}[htbp]
\centering
\tikzset{t1 hadron/.style={draw,solid,line width=1.5pt,line cap=butt}}
\begin{tikzpicture}[fd]
  \def\toneX{1.3}
  % heights: blob top 0.45, vertices 1.35, 2.5, 3.85, photon vertex 5.0; dots 2.95--3.25
  % left chain (amplitude)
  \draw[quark] (-\toneX,0.43) -- (-\toneX,1.35);
  \draw[quark] (-\toneX,1.35) -- (-\toneX,2.5);
  \draw[plain] (-\toneX,2.5) -- (-\toneX,2.82);
  \draw[line width=0.8pt,dash pattern=on 0pt off 3pt,line cap=round] (-\toneX,2.95) -- (-\toneX,3.25);
  \draw[quark] (-\toneX,3.38) -- (-\toneX,3.85);
  \draw[quark] (-\toneX,3.85) -- (-\toneX,5.0);
  % right chain (conjugate amplitude)
  \draw[quark] (\toneX,1.35) -- (\toneX,0.43);
  \draw[quark] (\toneX,2.5) -- (\toneX,1.35);
  \draw[plain] (\toneX,2.82) -- (\toneX,2.5);
  \draw[line width=0.8pt,dash pattern=on 0pt off 3pt,line cap=round] (\toneX,2.95) -- (\toneX,3.25);
  \draw[quark] (\toneX,3.85) -- (\toneX,3.38);
  \draw[quark] (\toneX,5.0) -- (\toneX,3.85);
  % struck quark and photons
  \draw[quark=0.25] (-\toneX,5.0) -- (\toneX,5.0);
  \draw[photon] (-2.35,5.75) -- (-\toneX,5.0);
  \draw[photon] (\toneX,5.0) -- (2.35,5.75);
  % rungs: real partons crossing the cut
  \foreach \y in {1.35,2.5,3.85} {
    \draw[gluon] (-\toneX,\y) -- (\toneX,\y);
    \node[vertex] at (-\toneX,\y) {}; \node[vertex] at (\toneX,\y) {};}
  \node[vertex] at (-\toneX,5.0) {}; \node[vertex] at (\toneX,5.0) {};
  % hadron blob (input PDF)
  \draw[t1 hadron] (-3.2,0) -- (3.2,0);
  \draw[hadron blob] (0,0) ellipse[x radius=1.85,y radius=0.45];
  \node at (-0.9,0) {$f(\xi_0;\mu_0)$};
  \draw[cut] (0,-0.75) -- (0,5.9);
  % labels: momentum fractions (left), splitting variables (vertices), k_T (right)
  \node[left] at (-\toneX,0.89) {$\xi_0$};
  \node[left] at (-\toneX,1.92) {$\xi_1=z_1\xi_0$};
  \node[left] at (-\toneX,3.6) {$\xi_{n-1}$};
  \node[left] at (-\toneX,4.42) {$\xi_n=z_n\xi_{n-1}\simeq x$};
  \foreach \y/\z in {1.35/z_1,2.5/z_2,3.85/z_n}
    \node[right,font=\footnotesize,inner sep=2pt] at (-\toneX,\y+0.31) {$\z$};
  \node[right] at (\toneX,0.89) {$k_{\mathrm T0}\sim\mu_0$};
  \node[right] at (\toneX,1.92) {$k_{\mathrm T1}$};
  \node[right] at (\toneX,3.6) {$k_{\mathrm T,n-1}$};
  \node[right] at (\toneX,4.42) {$k_{\mathrm Tn}$};
  \node[above left,inner sep=1pt] at (-2.35,5.75) {$\gamma^*(q)$};
  \node[left] at (-3.2,0) {$P$}; \node[right] at (3.2,0) {$P$};
  % ordering arrow
  \draw[lcaxis,line width=0.7pt] (3.75,0.89) -- (3.75,5.1);
  \node[right,align=left,font=\small] at (3.9,3.0)
    {$\mu_0^2\ll k_{\mathrm T1}^2\ll\cdots$\\[2pt]$\quad\ll k_{\mathrm T,n-1}^2\ll k_{\mathrm Tn}^2\ll Q^2$};
\end{tikzpicture}
\caption{DGLAP 演化在领头对数近似下的梯形图（物理规范中的非单态部分；单态部分的竖线也可以是胶子）。
自初始尺度 $\mu_0$ 的部分子分布 $f(\xi_0;\mu_0)$ 出发，第 $i$ 个分裂顶点把动量分数由 $\xi_{i-1}$ 变为 $\xi_i=z_i\xi_{i-1}$，
同时发射一个穿过末态割线的实部分子（带走 $(1-z_i)\xi_{i-1}$）。在横动量强序 $\mu_0^2\ll k_{\mathrm T1}^2\ll\cdots\ll k_{\mathrm Tn}^2\ll Q^2$ 的区域，
每一级给出 $\frac{\as}{2\pi}\int\frac{\dd k_{\mathrm T}^2}{k_{\mathrm T}^2}P(z_i)$，$n$ 级梯形图贡献 $\frac{1}{n!}\big(\frac{\as}{2\pi}\ln\frac{Q^2}{\mu_0^2}\big)^n$ 乘以 $n$ 重卷积 $P\otimes\cdots\otimes P$（固定耦合）；
连同提供加号分布的虚修正，对 $n$ 求和即得领头阶 DGLAP 方程的解。}
\label{fig:dglap-ladder}
\end{figure}
%%% END fig:dglap-ladder

\begin{proposition}[领头阶分裂函数 \tagR]\label{prop:pdf-lo-kernels}
（$R_{\mathrm{PT}}$：在重整化微扰论中对红外正规化的部分子态严格。）
在单圈阶（领头阶）：
(a) 由算符定义算出的 $D$ 维实辐射核（夸克自旋平均 $\frac12$，胶子极化平均 $\frac1{D-2}$）为
\begin{equation}\label{eq:pdf-Phat}
\begin{aligned}
\hat P_{qq}&=C_F\Big[\frac{1+z^2}{1-z}-\varepsilon(1-z)\Big], &
\hat P_{gq}&=C_F\Big[\frac{1+(1-z)^2}{z}-\varepsilon z\Big],\\
\hat P_{qg}&=T_F\Big[1-\frac{2z(1-z)}{1-\varepsilon}\Big], &
\hat P_{gg}&=2C_A\Big[\frac{z}{1-z}+\frac{1-z}{z}+z(1-z)\Big];
\end{aligned}
\end{equation}
% VF-PDF-6
(b) 领头阶分裂函数为
\begin{equation}\label{eq:pdf-Pij}
\begin{gathered}
P_{qq}=C_F\Big[\frac{1+z^2}{1-z}\Big]_+=C_F\Big[\frac{1+z^2}{(1-z)_+}+\frac32\delta(1-z)\Big],\qquad P_{qg}=T_F\big[z^2+(1-z)^2\big],\\
P_{gq}=C_F\frac{1+(1-z)^2}{z},\qquad P_{gg}=2C_A\Big[\frac{z}{(1-z)_+}+\frac{1-z}{z}+z(1-z)\Big]+\frac{\beta_0}{2}\delta(1-z),
\end{gathered}
\end{equation}
% VF-PDF-7
$\frac{z}{(1-z)_+}\equiv z[\frac1{1-z}]_+$；单态演化 $\frac{\dd}{\dd\ln\mu^2}\binom{\Sigma}{g}=\frac{\as}{2\pi}\begin{pmatrix}P_{qq}&2n_fP_{qg}\\P_{gq}&P_{gg}\end{pmatrix}\otimes\binom{\Sigma}{g}$，$\Sigma=\sum_k(q_k+\bar q_k)$；求和规则 $\int_0^1P_{qq}=0$，$\int_0^1z[P_{qq}+P_{gq}]=0$，$\int_0^1z[P_{gg}+2n_fP_{qg}]=0$；
(c) Mellin 矩 $\gamma_{ij}(N)=\int_0^1z^{N-1}P_{ij}$ 为
\begin{equation}\label{eq:pdf-gammas}
\begin{gathered}
\gamma_{qq}=C_F\Big[\frac32+\frac{1}{N(N+1)}-2S_1(N)\Big],\qquad \gamma_{qg}=T_F\frac{N^2+N+2}{N(N+1)(N+2)},\\
\gamma_{gq}=C_F\frac{N^2+N+2}{(N-1)N(N+1)},\qquad \gamma_{gg}=2C_A\Big[\frac{1}{N(N-1)}+\frac{1}{(N+1)(N+2)}-S_1(N)\Big]+\frac{\beta_0}{2};
\end{gathered}
\end{equation}
% VF-PDF-8
(d) $N=2$ 单态矩阵为 $\begin{pmatrix}-\frac43C_F&\frac23n_fT_F\\\frac43C_F&-\frac23n_fT_F\end{pmatrix}$，本征值 $0$ 与 $-\frac{4C_F+2n_fT_F}{3}$；对 SU(3)，$\langle x\rangle_\Sigma\to\frac{3n_f}{16+3n_f}$；非单态矩的领头阶演化为 $M^{\mathrm{NS}}_n(Q^2)=M^{\mathrm{NS}}_n(Q_0^2)[\as(Q^2)/\as(Q_0^2)]^{-2\gamma_{qq}(n)/\beta_0}$；
% VF-PDF-9
(e) 与文献约定的换算：MVV（$a_s=\frac{\as}{4\pi}$，$\gamma^{(n)}=-\int x^{N-1}P^{(n)}$）有 $P^{(0)}_{\mathrm{MVV}}=2P$、$\gamma^{(0)}_{\mathrm{MVV}}=-2\gamma$；Gross--Wilczek 与 Georgi--Politzer 的 $\mu\frac{\dd O_N}{\dd\mu}=-\gamma_OO_N$ 有 $\gamma_O=-\frac{\as}{\pi}\gamma(N)$；定域算符按 $\frac{\dd O}{\dd\ln\mu^2}=\frac{\as}{2\pi}\Gamma O$ 有 $\Gamma_{qq}=\gamma_{qq}$、$\Gamma_{q_kg}=2\gamma_{qg}$、$\Gamma_{gq}=\gamma_{gq}$、$\Gamma_{gg}=\gamma_{gg}$。
% VF-PDF-10
\uses{thm:pdf-dglap,def:pdf-quark,def:pdf-gluon,thm:pdf-moments-local}
\end{proposition}

\begin{proof}
(a) 取 $A^+=0$ 规范（$W=1$），靶为无质量在壳部分子 $p=(p^+,0^-,\bm 0_{\mathrm T})$，末态为一个无质量实部分子 $l=\big((1-z)p^+,\frac{\bm l_{\mathrm T}^2}{2(1-z)p^+},\bm l_{\mathrm T}\big)$，$k=p-l$，$k^+=zp^+$，$k^2=-\bm l_{\mathrm T}^2/(1-z)$。以 $q\leftarrow q$ 为例：振幅 $\langle l,a,\epsilon|\psi(0)|p\rangle=-g\,t^a\frac{\sla k\sla\epsilon^*}{k^2}u(p)$，物理极化和 $d^{\mu\nu}(l)=-g^{\mu\nu}+\frac{l^\mu n_B^\nu+n_B^\mu l^\nu}{l^+}$。代入\cref{eq:pdf-cut}（含其中的 $\frac12$），对靶自旋平均、对色求和平均（$C_F$），一粒子相空间中的 $l^+$ 被 $\delta$ 函数固定，得
\[
f^{\mathrm{real}}_{q/q}(z)=\frac{g^2C_F}{(2\pi)^{D-1}}\!\int\!\dd^{D-2}\bm l_{\mathrm T}\,
\frac{\TrD[\sla p\gamma_\nu\sla k\gamma^+\sla k\gamma_\mu]d^{\mu\nu}(l)}{4\cdot2l^+(k^2)^2}
=\frac{g^2}{(2\pi)^{D-1}}\!\int\!\dd^{D-2}\bm l_{\mathrm T}\,\frac{\hat P_{qq}(z;\varepsilon)}{\bm l_{\mathrm T}^2},
\]
末步是 $D$ 维 Dirac 迹的直接计算。其余三道相同，只需：$g\leftarrow q$ 用胶子算符顶点 $X^{j\beta}=(k^+g^{j\alpha}-k^jn_B^\alpha)d_\alpha{}^\beta(k)$ 与\cref{eq:pdf-fg}的 $1/(zp^+)$；$q\leftarrow g$ 用 $\sum v\bar v=\sla l$ 与初态极化平均 $-g_\perp^{\mu\nu}/(D-2)$；$g\leftarrow g$ 用三胶子顶点与两个物理极化和。四道结果作为 $D$ 的有理函数精确等于\cref{eq:pdf-Phat}。
(b) 有红外正规化时 $\int\dd^{2-2\varepsilon}\bm l_{\mathrm T}/(\bm l_{\mathrm T}^2+M^2)=\pi^{1-\varepsilon}\Gamma(\varepsilon)M^{-2\varepsilon}$，而 $\MSbar$ 中 $\frac{g_0^2}{(2\pi)^{3-2\varepsilon}}\pi^{1-\varepsilon}=\frac{\as}{2\pi}\mu^{2\varepsilon}\ee^{\varepsilon\gamma_E}$，故实辐射的紫外极点为 $\frac{\as}{2\pi}\frac1\varepsilon\hat P_{ij}(z;0)$，对 $z<1$ 给出 $Z^{(1)}_{ij}=-\frac{\as}{2\pi}\hat P_{ij}(z;0)$，由\cref{thm:pdf-dglap}，$P_{ij}=\hat P_{ij}(z;0)$（$z<1$）。$z=1$ 处只能出现 $\delta(1-z)$ 与正规化 $1/(1-z)$ 的加号分布，且只有对角道有虚图。$\delta$ 项由守恒律确定：夸克数流 $\bar\psi\gamma^\mu\psi$ 守恒，其荷是好量子数，故 $\bar\psi\gamma^+\psi$ 不需重整化，由\cref{thm:pdf-moments-local}（$n=1$）与\cref{thm:pdf-dglap}得 $\int_0^1P_{qq}=0$，即 $b_q=C_F\int_0^1(1+z)\dd z=\frac32C_F$；能动张量守恒且其荷为四动量，故 $T^{++}=\sum_qO^{(2)}_q+O^{(2)}_g$ 不需重整化（\cite{Collins1984}），由 $n=2$ 得 $\sum_i\int_0^1zP_{ij}=0$。胶子列用 $\int_0^1\frac{z^2}{(1-z)_+}=-\frac32$、$\int_0^1(1-z)=\frac12$、$\int_0^1z^2(1-z)=\frac1{12}$、$2n_fT_F\int_0^1z[z^2+(1-z)^2]=\frac23n_fT_F$ 得 $b_g=\frac{11}6C_A-\frac23n_fT_F=\frac{\beta_0}2$；夸克列 $\int_0^1zP_{qq}=-\frac43C_F$ 与 $\int_0^1zP_{gq}=\frac43C_F$ 之和为零，是不含待定参数的独立检验。两种写法的 $P_{qq}$ 相差 $\delta(1-z)\int_0^1\frac{2-(1+z^2)}{1-z}\dd z=\frac32\delta(1-z)$。电荷共轭给出 $P_{\bar q\bar q}=P_{qq}$ 等，领头阶 $P_{q\bar q}=0$；单态矩阵中的 $2n_fP_{qg}$ 来自每味夸克与反夸克各一份。
(c) 由 $\int_0^1\frac{z^{N-1}-1}{1-z}\dd z=-S_1(N-1)$（$S_1(N)=\psi(N+1)+\gamma_E$）得 $\int_0^1z^{N-1}[\frac{1+z^2}{1-z}]_+=\frac32-S_1(N-1)-S_1(N+1)=\frac32+\frac{1}{N(N+1)}-2S_1(N)$；其余直接积分。
(d) $\gamma_{qq}(2)=-\frac43C_F$，$\gamma_{gq}(2)=\frac43C_F$，$2n_f\gamma_{qg}(2)=\frac23n_fT_F$，$\gamma_{gg}(2)=-\frac23n_fT_F$；两列之和为零，故 $(1,1)$ 是本征值 $0$ 的左本征矢，另一本征值等于迹；本征值 $0$ 的右本征矢 $\propto(2n_fT_F,4C_F)$ 给出 $\langle x\rangle_\Sigma\to\frac{2n_fT_F}{4C_F+2n_fT_F}=\frac{3n_f}{16+3n_f}$。非单态矩满足 $\frac{\dd M_n}{\dd\ln\mu^2}=\frac{\as}{2\pi}\gamma_{qq}(n)M_n$，与 $\frac{\dd\as}{\dd\ln\mu^2}=-\frac{\beta_0}{4\pi}\as^2$ 联立积分即得幂律。
(e) 由各文献的定义直接推出。
\end{proof}

$\delta$ 项也可直接由虚图得到，两条路一致：$qq$ 道见\cref{prop:pdf-pqq-operator}（$\frac32=2-\frac12$），$gg$ 道的 $\frac{\beta_0}2$ 见\cref{prop:pdf-local-anomalous}。领头阶非单态时，MVV 的 $\gamma^{(0)}_{\mathrm{ns}}=C_F[4S_1(N)-3-\frac{2}{N(N+1)}]$\cite{MochVermaserenVogt2004}、Gross--Wilczek 与 Georgi--Politzer 的 $\gamma_O=\frac{\as}{2\pi}C_F[1-\frac{2}{N(N+1)}+4\sum_{j=2}^N\frac1j]$\cite{GrossWilczek1974,GeorgiPolitzer1974}，都等于 $\gamma_{qq}(N)$ 乘以 (e) 中的因子；延拓到非整数 $N$ 的唯一性见\cref{prop:appK-continuation}。四种分裂见\cref{fig:splitting-vertices}。

% -----------------------------------------------------------------------------------------------------
\section{单圈计算：由算符定义得到 $P_{qq}$}\label{sec:pdf-oneloop}

本节在 Feynman 规范下直接由\cref{eq:pdf-fq}算出夸克中夸克分布的单圈值；Wilson 线显式出现，每个图的紫外极点可以单独读出。

\begin{proposition}[Feynman 规范下由算符定义计算 $P_{qq}$ \tagR]\label{prop:pdf-pqq-operator}
（$R_{\mathrm{PT}}$：在重整化微扰论中对红外正规化的部分子态严格。）
在单圈阶，取 Feynman 规范，以胶子质量 $m$ 作红外正规化、维数正规化作紫外正规化，靶为自旋与色平均的在壳无质量夸克 $p=(p^+,0^-,\bm 0_{\mathrm T})$，记
\begin{equation}\label{eq:pdf-Aeps}
A_\varepsilon\equiv\frac{\as}{2\pi}C_F\,\ee^{\varepsilon\gamma_E}\Gamma(\varepsilon)\Big(\frac{\mu^2}{m^2}\Big)^{\varepsilon}=\frac{\as}{2\pi}C_F\Big[\frac1\varepsilon+\ln\frac{\mu^2}{m^2}+\order{\varepsilon}\Big].
\end{equation}
(a) 实辐射部分为 $f_R(\xi)=A_\varepsilon\big[(1-\varepsilon)^2(1-\xi)\xi^{-\varepsilon}+\frac{2\xi^{1-\varepsilon}}{1-\xi}\big]$。
(b) Wilson 线的实图与虚图只在组合中构成分布 $A_\varepsilon[\frac{2\xi^{1-\varepsilon}}{1-\xi}]_+$（$\xi\to1$ 的端点奇异性不受维数正规化调控）；外线自能贡献 $-A_\varepsilon\frac{1-\varepsilon}{2-\varepsilon}\delta(1-\xi)$。
(c) 裸单圈分布对检验函数 $h$ 为
\begin{equation}\label{eq:pdf-bare1}
\langle f^{(1)}_{\mathrm{bare}},h\rangle=A_\varepsilon\Big\{\int_0^1\!\dd\xi\Big[(1-\varepsilon)^2(1-\xi)\xi^{-\varepsilon}h(\xi)+\frac{2\xi^{1-\varepsilon}}{1-\xi}\big(h(\xi)-h(1)\big)\Big]-\frac{1-\varepsilon}{2-\varepsilon}h(1)\Big\},
\end{equation}
夸克数对任意 $\varepsilon$ 精确守恒：$\int_0^1f^{(1)}_{\mathrm{bare}}=0$。
(d) 紫外极点逐图（单位 $\frac{\as}{2\pi}\frac{C_F}{\varepsilon}$）：梯形图 $(1-\xi)$，Wilson 线图 $[\frac{2\xi}{1-\xi}]_+$，自能 $-\frac12\delta(1-\xi)$；其和为 $P_{qq}/C_F$。
(e) $\MSbar$ 重整化后（依赖于红外正规化，不是物理量）
\begin{equation}\label{eq:pdf-ren1}
f^{(1)}_{q/q}(\xi;\mu)=\frac{\as}{2\pi}\Big\{P_{qq}(\xi)\ln\frac{\mu^2}{m^2}+C_F\Big[-(1-\xi)(2+\ln\xi)-\Big[\frac{2\xi\ln\xi}{1-\xi}\Big]_++\frac14\delta(1-\xi)\Big]\Big\}.
\end{equation}
(f) 在光锥规范中只有梯形图，它单独给出同一 $D$ 维核 $\hat P_{qq}$。
（本命题只涉及部分子态。）
% VF-PDF-3
\uses{def:pdf-quark,def:prelim-feynman-rules,prop:prelim-wilson-rules,def:prelim-msbar}
\end{proposition}

\begin{proof}
\emph{图。}用割图表示\cref{eq:pdf-cut}。$\order{\as}$ 的图见\cref{fig:pdf-oneloop}：(a) 梯形图 R1；(b) 实胶子一端在夸克线、另一端在割线另一侧的 Wilson 线上（R2 及镜像）；(c) 虚图：夸克线—Wilson 线顶点修正 V1 与外线自能 V2（及镜像）；(d) 两端都在 Wilson 线上的图 $\propto n_B^2=0$。
\emph{实辐射。}由\cref{prop:prelim-wilson-rules}，$W[\infty n_B,0]$ 吸收流入动量 $-l$（即发射出射胶子 $l$）的单胶子项为 $g\,t^an_B^\mu/(-n_B\cdot l-\ii0)$，故
\begin{equation}\label{eq:pdf-amp}
\langle l,a,\epsilon|W[\infty n_B,0]\psi(0)|p\rangle=-g\,t^a\epsilon^*_\mu\Big[\frac{\sla k\gamma^\mu}{k^2}+\frac{n_B^\mu}{n_B\cdot l+\ii0}\Big]u(p),\qquad k=p-l .
\end{equation}
相对符号由规范不变性确定：$\sla lu(p)=-\sla ku(p)$，故 $l_\mu[\cdots]u=(\sla k\sla l/k^2+1)u=0$；因此有质量胶子极化和中的 $l^\mu l^\nu/m^2$ 项不贡献，可用 $-g_{\mu\nu}$。运动学：$l=\big((1-\xi)p^+,\frac{\bm l_{\mathrm T}^2+m^2}{2(1-\xi)p^+},\bm l_{\mathrm T}\big)$，$n_B\cdot k=\xi p^+$，$k^2=-(\bm l_{\mathrm T}^2+\xi m^2)/(1-\xi)$。对靶自旋平均后，$D$ 维迹给出（$\Gamma^\mu$ 为\cref{eq:pdf-amp}方括号中的矩阵）
\[
\tfrac12\TrD[\sla p\,\bar\Gamma^\nu\gamma^+\Gamma^\mu](-g_{\mu\nu})g^2C_F=g^2C_Fp^+\Big[\underbrace{\frac{4(1-\varepsilon)(1-\xi)^2\bm l_{\mathrm T}^2}{(\bm l_{\mathrm T}^2+\xi m^2)^2}}_{\text{R1}}+\underbrace{\frac{8\xi}{\bm l_{\mathrm T}^2+\xi m^2}}_{\text{R2+镜像}}\Big],
\]
R1 来自 $\gamma^\mu\sla k\gamma^+\sla k\gamma_\mu=-(D-2)\sla k\gamma^+\sla k$，R2 来自 $\gamma^+\sla k\gamma^+=2k^+\gamma^+$。一粒子相空间中 $l^+$ 被 $\delta$ 函数固定，连同\cref{eq:pdf-cut}中的 $\frac12$，
\[
f_R(\xi)=\frac{g^2C_F}{(2\pi)^{D-1}}\int\dd^{D-2}\bm l_{\mathrm T}\Big[\frac{(1-\varepsilon)(1-\xi)\bm l_{\mathrm T}^2}{(\bm l_{\mathrm T}^2+\xi m^2)^2}+\frac{2\xi}{(1-\xi)(\bm l_{\mathrm T}^2+\xi m^2)}\Big].
\]
用 $\int\dd^{d}\bm l\,\frac{(\bm l^2)^a}{(\bm l^2+M^2)^b}=\pi^{d/2}\frac{\Gamma(a+d/2)\Gamma(b-a-d/2)}{\Gamma(d/2)\Gamma(b)}(M^2)^{a+d/2-b}$（$d=2-2\varepsilon$，$M^2=\xi m^2$）与 $\MSbar$ 中的 $\frac{g_0^2}{(2\pi)^{3-2\varepsilon}}\pi^{1-\varepsilon}=\frac{\as}{2\pi}\mu^{2\varepsilon}\ee^{\varepsilon\gamma_E}$ 即得 (a)。
\emph{虚 Wilson 线图 V1。}用 $\bar u\gamma^+(\sla p-\sla l)\gamma^+u=2(p^+-l^+)\bar u\gamma^+u$，
\[
\frac{V_1}{\text{树图}}=\ii g^2C_F\int\frac{\dd^Dl}{(2\pi)^D}\,\frac{2(p^+-l^+)}{(l^+-\ii0)(l^2-m^2+\ii0)((p-l)^2+\ii0)} .
\]
两个传播子在 $l^-$ 平面的极点 $\frac{\bm l_{\mathrm T}^2+m^2-\ii0}{2l^+}$ 与 $-\frac{\bm l_{\mathrm T}^2-\ii0}{2(p^+-l^+)}$ 只当 $0<l^+<p^+$ 时位于实轴两侧；闭合围道得 $\frac{\pi\ii\,\theta(0<l^+<p^+)}{p^+(\bm l_{\mathrm T}^2+xm^2)}$，$l^+=(1-x)p^+$。用 $\frac{2(p^+-l^+)}{l^+p^+}\dd l^+=\frac{2x}{1-x}\dd x$ 与 $\ii\cdot\frac{\pi\ii}{(2\pi)^D}=-\frac{1}{2(2\pi)^{D-1}}$，加上镜像，
\[
\frac{V_1+V_1^\dagger}{\text{树图}}=-\frac{g^2C_F}{(2\pi)^{D-1}}\int_0^1\!\dd x\!\int\!\dd^{D-2}\bm l_{\mathrm T}\frac{2x}{(1-x)(\bm l_{\mathrm T}^2+xm^2)}=-A_\varepsilon\int_0^1\!\dd x\,\frac{2x^{1-\varepsilon}}{1-x},
\]
它乘以 $\delta(1-\xi)$，被积函数恰为 R2 在 $\xi=x$ 处的相反数，所以 R2 与 V1 合成 $A_\varepsilon[2\xi^{1-\varepsilon}/(1-\xi)]_+$；端点 $x\to1$（$l^+\to0$）的发散只在这一组合中相消，必须在固定 $x=\xi$ 处先组合再积分。
\emph{外线自能 V2。}$\Sigma(p)=\sla p\,\Sigma_V(p^2)$，$\Sigma_V(0)=(2-D)\frac{g_0^2C_F}{(4\pi)^{D/2}}\Gamma(\varepsilon)\int_0^1\dd y\,(1-y)[(1-y)m^2]^{-\varepsilon}=-A_\varepsilon\frac{1-\varepsilon}{2-\varepsilon}$，其极点 $-\frac{\as}{4\pi}\frac{C_F}\varepsilon$ 即 Feynman 规范的 $\delta Z_2$；两条外线的 LSZ 留数合计贡献 $(R-1)\delta(1-\xi)=\Sigma_V(0)\delta(1-\xi)+\order{\as^2}$。这证明了 (b)，相加得 (c)。
\emph{夸克数、极点、有限部分。}$(1-\varepsilon)^2\int_0^1(1-\xi)\xi^{-\varepsilon}\dd\xi=\frac{1-\varepsilon}{2-\varepsilon}$，加号分布积分为零，故 $h=1$ 时\cref{eq:pdf-bare1}为零。由 $\ee^{\varepsilon\gamma_E}\Gamma(\varepsilon)=\frac1\varepsilon+\order\varepsilon$（无常数项），极点为 $\frac{\as}{2\pi}\frac{C_F}\varepsilon\{(1-\xi)+[\frac{2\xi}{1-\xi}]_+-\frac12\delta\}$，而用 $[\frac{2\xi}{1-\xi}]_+=\frac{2\xi}{(1-\xi)_+}+2\delta(1-\xi)$，花括号等于 $\frac{1+\xi^2}{(1-\xi)_+}+\frac32\delta=[\frac{1+\xi^2}{1-\xi}]_+$，即 (d)。减去 $\frac{\as}{2\pi}\frac1\varepsilon P_{qq}$ 后，$\varepsilon^0$ 项是 $\frac{\as}{2\pi}P_{qq}\ln\frac{\mu^2}{m^2}$ 加上花括号对 $\varepsilon$ 的一阶导数：$\frac{\dd}{\dd\varepsilon}[(1-\varepsilon)^2(1-\xi)\xi^{-\varepsilon}]_0=-(1-\xi)(2+\ln\xi)$，$\frac{\dd}{\dd\varepsilon}[2\xi^{1-\varepsilon}]_0=-2\xi\ln\xi$，$-\frac{\dd}{\dd\varepsilon}[\frac{1-\varepsilon}{2-\varepsilon}]_0=\frac14$，即 (e)。
\emph{光锥规范。}取 $d^{\mu\nu}(l)=-g^{\mu\nu}+(l^\mu n_B^\nu+n_B^\mu l^\nu)/(n_B\cdot l)$，$W=1$，只剩梯形图，其迹乘以 $\bm l_{\mathrm T}^2$（$m=0$）同样等于 $\frac{1+\xi^2}{1-\xi}-\varepsilon(1-\xi)$，即 (f)：Feynman 规范中由 Wilson 线承担的部分，在光锥规范中由极化和的 $1/(n_B\cdot l)$ 项承担。
\end{proof}

%%% BEGIN fig:pdf-oneloop
\begin{figure}[htbp]
\centering
% frame common to all graphs: cut, Wilson line (arrow at position #1), operator vertices, external labels
\newcommand{\toneW}[1]{%
  \draw[cut] (0,-0.22) -- (0,2.6);
  \draw[wilson=#1] (-0.75,1.9) -- (0.75,1.9);
  \node[opvertex] at (-0.75,1.9) {}; \node[opvertex] at (0.75,1.9) {};
  \node[below] at (-0.75,0) {$p$}; \node[below] at (0.75,0) {$p$};}
\begin{subfigure}[b]{0.24\textwidth}\centering
\begin{tikzpicture}[fd]
  \draw[quark] (-0.75,0) -- (-0.75,0.9);
  \draw[quark] (-0.75,0.9) -- (-0.75,1.9);
  \draw[quark] (0.75,1.9) -- (0.75,0.9);
  \draw[quark] (0.75,0.9) -- (0.75,0);
  \draw[gluon] (-0.75,0.9) -- (0.75,0.9);
  \toneW{0.25}
  \node[vertex] at (-0.75,0.9) {}; \node[vertex] at (0.75,0.9) {};
  \path[mom'={$l$},mom pos=0.76,mom length=0.42cm,mom sep=0.27cm] (-0.75,0.9) -- (0.75,0.9);
  \path[mom={$k$},mom length=0.45cm] (-0.75,0.9) -- (-0.75,1.9);
\end{tikzpicture}
\caption{实胶子发射}
\end{subfigure}\hspace{2em}
\begin{subfigure}[b]{0.56\textwidth}\centering
\begin{tikzpicture}[fd]
  \draw[quark] (-0.75,0) -- (-0.75,0.8);
  \draw[quark] (-0.75,0.8) -- (-0.75,1.9);
  \draw[quark] (0.75,1.9) -- (0.75,0);
  \draw[gluon] (-0.75,0.8) to[out=10,in=-100] (0.4,1.9);
  \toneW{0.25}
  \node[vertex] at (-0.75,0.8) {}; \node[vertex] at (0.4,1.9) {};
\end{tikzpicture}\hspace{1.5em}
\begin{tikzpicture}[fd]
  \draw[quark] (-0.75,0) -- (-0.75,1.9);
  \draw[quark] (0.75,1.9) -- (0.75,0.8);
  \draw[quark] (0.75,0.8) -- (0.75,0);
  \draw[gluon] (-0.4,1.9) to[out=-80,in=170] (0.75,0.8);
  \toneW{0.75}
  \node[vertex] at (0.75,0.8) {}; \node[vertex] at (-0.4,1.9) {};
\end{tikzpicture}
\caption{夸克线与另一侧 Wilson~线间的实胶子（及镜像）}
\end{subfigure}

\bigskip
\begin{subfigure}[b]{0.68\textwidth}\centering
\begin{tikzpicture}[fd]
  \draw[quark] (-0.75,0) -- (-0.75,0.8);
  \draw[quark] (-0.75,0.8) -- (-0.75,1.9);
  \draw[quark] (0.75,1.9) -- (0.75,0);
  \draw[gluon] (-0.75,0.8) to[out=0,in=-90] (-0.25,1.9);
  \toneW{0.75}
  \node[vertex] at (-0.75,0.8) {}; \node[vertex] at (-0.25,1.9) {};
\end{tikzpicture}\hspace{0.9em}
\begin{tikzpicture}[fd]
  \draw[quark] (-0.75,0) -- (-0.75,1.9);
  \draw[quark] (0.75,1.9) -- (0.75,0.8);
  \draw[quark] (0.75,0.8) -- (0.75,0);
  \draw[gluon] (0.25,1.9) to[out=-90,in=180] (0.75,0.8);
  \toneW{0.25}
  \node[vertex] at (0.75,0.8) {}; \node[vertex] at (0.25,1.9) {};
\end{tikzpicture}\hspace{0.9em}
\begin{tikzpicture}[fd]
  \draw[quark] (-0.75,0) -- (-0.75,0.45);
  \draw[quark] (-0.75,0.45) -- (-0.75,1.3);
  \draw[quark] (-0.75,1.3) -- (-0.75,1.9);
  \draw[quark] (0.75,1.9) -- (0.75,0);
  \draw[gluon] (-0.75,0.45) arc[start angle=-90,end angle=-270,radius=0.425];
  \toneW{0.25}
  \node[vertex] at (-0.75,0.45) {}; \node[vertex] at (-0.75,1.3) {};
\end{tikzpicture}\hspace{0.9em}
\begin{tikzpicture}[fd]
  \draw[quark] (-0.75,0) -- (-0.75,1.9);
  \draw[quark] (0.75,1.9) -- (0.75,1.3);
  \draw[quark] (0.75,1.3) -- (0.75,0.45);
  \draw[quark] (0.75,0.45) -- (0.75,0);
  \draw[gluon] (0.75,1.3) arc[start angle=90,end angle=-90,radius=0.425];
  \toneW{0.25}
  \node[vertex] at (0.75,0.45) {}; \node[vertex] at (0.75,1.3) {};
\end{tikzpicture}
\caption{不穿过割线的虚图（两两互为镜像）}
\end{subfigure}\hfill
\begin{subfigure}[b]{0.28\textwidth}\centering
\begin{tikzpicture}[fd]
  \draw[quark] (-0.75,0) -- (-0.75,1.9);
  \draw[quark] (0.75,1.9) -- (0.75,0);
  \draw[gluon] (-0.5,1.9) arc[start angle=180,end angle=0,radius=0.5];
  \toneW{0.36}
  \node[vertex] at (-0.5,1.9) {}; \node[vertex] at (0.5,1.9) {};
  \node[font=\footnotesize] at (0,-0.68) {$\propto n^2=0$};
\end{tikzpicture}
\caption{Wilson~线之间的胶子}
\end{subfigure}
\caption{Feynman 规范下夸克中夸克分布 $f_{q/q}(\xi)$ 的 $\order{\as}$ 割图（靶为在壳夸克 $p$，$k^+=\xi p^+$；记号同\cref{fig:pdf-operator}）。
(a) 实胶子发射（梯形图），给出 $P_{qq}$ 中的 $C_F(1-\xi)$；(b) 实胶子由夸克线连到割线另一侧的 Wilson 线，及其镜像；
(c) 不穿过割线的虚图：夸克线—Wilson 线顶点修正（前两图）与外线自能（后两图）。
(b) 与 (c) 中夸克线—Wilson 线图在 $\xi\to1$ 的端点奇异性只在实、虚之和中相消，合成 $C_F[2\xi/(1-\xi)]_+$；
外线自能的紫外发散（夸克场重整化）贡献 $-\tfrac12C_F\delta(1-\xi)$；总和为 $P_{qq}=C_F[(1+\xi^2)/(1-\xi)]_+$。
(d) 连接割线两侧 Wilson 线的胶子正比于 $n^2=0$，在 Feynman 规范中为零（同侧的 Wilson 线自能亦然）。}
\label{fig:pdf-oneloop}
\end{figure}
%%% END fig:pdf-oneloop

\begin{example}[单圈夸克 PDF（逐图）]\label{ex:pdf-oneloop}
按\cref{fig:pdf-oneloop}逐图：R1 给出 $A_\varepsilon(1-\varepsilon)^2(1-\xi)\xi^{-\varepsilon}$，R2 与 V1 给出 $A_\varepsilon[\frac{2\xi^{1-\varepsilon}}{1-\xi}]_+$，V2 给出 $-A_\varepsilon\frac{1-\varepsilon}{2-\varepsilon}\delta(1-\xi)$，(d) 为零；$\delta$ 项 $\frac32=2-\frac12$ 中，$2$ 来自把 R2 与 V1 的组合改写为标准形式，$-\frac12$ 来自夸克场重整化，梯形图不含 $\delta$ 项。数值上，取 $\mu=m$、$\xi=\frac12$，\cref{eq:pdf-ren1}的正规部分为 $\frac{\as C_F}{2\pi}[-\frac12(2-\ln2)+2\ln2]=\frac{\as C_F}{2\pi}(\frac52\ln2-1)=0.7329\,\frac{\as C_F}{2\pi}$；$\ln(\mu^2/m^2)$ 每改变一个单位，该点的值改变 $\frac{\as}{2\pi}P_{qq}(\frac12)=\frac52\frac{\as C_F}{2\pi}$，正是 DGLAP 所要求的 $\mu$ 依赖。夸克数检验：$\int_0^1[-(1-\xi)(2+\ln\xi)]\dd\xi=-\frac14$，恰与 $\frac14\delta(1-\xi)$ 相消。有限部分依赖于红外正规化（纯维数正规化下为零，\cref{prop:pdf-partonic-dr}），只有紫外极点与 $\ln\mu^2$ 的系数与正规化无关。
\uses{prop:pdf-pqq-operator}
\end{example}

\begin{proposition}[纯维数正规化中的部分子分布 \tagR]\label{prop:pdf-partonic-dr}
（$R_{\mathrm{PT}}$：在重整化微扰论中对红外正规化的部分子态严格。）
设靶为无质量在壳部分子 $j$，维数正规化同时正规化紫外与红外发散。则裸部分子分布在所有阶都是 $f^{\mathrm{bare}}_{i/j}(\xi)=\delta_{ij}\delta(1-\xi)$；$\MSbar$ 重整化后
\begin{equation}\label{eq:pdf-partonicDR}
f_{i/j}(\xi;\mu)=Z_{ij}(\xi)=\delta_{ij}\delta(1-\xi)-\frac{\as}{2\pi}\frac1\varepsilon P_{ij}(\xi)+\order{\as^2},
\end{equation}
% VF-PDF-5
剩下的 $1/\varepsilon$ 是红外极点：$\int\frac{\dd^{2-2\varepsilon}\bm l_{\mathrm T}}{\bm l_{\mathrm T}^2}=\pi^{1-\varepsilon}\Gamma(1+\varepsilon)M^{-2\varepsilon}\big(\frac1{\varepsilon_{\mathrm{UV}}}-\frac1{\varepsilon_{\mathrm{IR}}}\big)=0$，$\MSbar$ 抵消项去掉 $1/\varepsilon_{\mathrm{UV}}$ 后留下 $-1/\varepsilon_{\mathrm{IR}}$。（本命题只涉及部分子态。）
\uses{def:pdf-quark,prop:pdf-invariance,lem:prelim-scaleless,ex:prelim-scaleless,prop:pdf-lo-kernels}
\end{proposition}

\begin{proof}
设 $F_L(\xi,p^+;\varepsilon)$ 是 $L\ge1$ 圈的裸贡献。它无量纲，唯一有量纲的变量是 $p^+$（$p^2=0$，无质量）；按量纲分析，把所有动量放大 $\lambda$ 倍时 $L$ 圈测度与 $g_0^{2L}\propto\mu_0^{2L\varepsilon}$ 给出 $F_L(\xi,\lambda p^+)=\lambda^{-2L\varepsilon}F_L(\xi,p^+)$；另一方面推动不变性（\cref{prop:pdf-invariance}(b)，对裸分布同样成立）给出 $F_L(\xi,\lambda p^+)=F_L(\xi,p^+)$。故 $(\lambda^{-2L\varepsilon}-1)F_L=0$ 对一切 $\lambda>0$ 成立，$\varepsilon\ne0$ 时 $F_L=0$——这正是\cref{lem:prelim-scaleless}的情形。所以 $f^{\mathrm{bare}}_{i/j}=\delta_{ij}\delta(1-\xi)$（\cref{prop:pdf-invariance}(c)），$f^{\mathrm{ren}}_{i/j}=\sum_kZ_{ik}\otimes\delta_{kj}\delta=Z_{ij}$，单圈的 $Z$ 由\cref{eq:pdf-Z}与\cref{prop:pdf-lo-kernels}给出\cref{eq:pdf-partonicDR}。单圈实辐射积分正是 $\int\dd^{2-2\varepsilon}\bm l_{\mathrm T}/\bm l_{\mathrm T}^2$：在任意 $M$ 处分开，$|\bm l_{\mathrm T}|>M$ 部分只在 $\varepsilon>0$ 时收敛，给出 $+1/\varepsilon_{\mathrm{UV}}$，$|\bm l_{\mathrm T}|<M$ 部分只在 $\varepsilon<0$ 时收敛，给出 $-1/\varepsilon_{\mathrm{IR}}$，各自延拓后相消（\cref{ex:prelim-scaleless}）。$\MSbar$ 抵消项与红外正规化无关，去掉 $1/\varepsilon_{\mathrm{UV}}$ 后剩下的 $-\frac{\as}{2\pi}\frac1{\varepsilon_{\mathrm{IR}}}P_{ij}$ 是部分子靶的共线发散。
\end{proof}

这正是把部分子截面变成系数函数的那一步：对部分子靶，因子化公式中的 PDF 是纯红外极点，所以系数函数等于“$\MSbar$ 减除了共线极点的部分子截面”（第~10~章）；忘记\cref{eq:pdf-partonicDR}右边的极点是红外的，共线减除中就很容易弄错符号。

\begin{proposition}[定域算符的单圈反常维度（含混合）\tagR]\label{prop:pdf-local-anomalous}
（$R_{\mathrm{PT}}$：在重整化微扰论中对红外正规化的部分子态严格。）
在单圈阶，对在壳、物理极化的无质量部分子态 $j$ 与扭度 2 定域算符 $O_q^{(N)}=\bar\psi\sla n_B(\ii n_B\cdot D)^{N-1}\psi$（单一味道）、$O_g^{(N)}=\sum_jG^{+j}(\ii D^+)^{N-2}G^{+j}$，定义 $O_i^R=\sum_jZ_{ij}O_j$，$Z=\mathbb 1-\frac{\as}{2\pi}\frac1\varepsilon\Gamma$（于是 $\frac{\dd O^R}{\dd\ln\mu^2}=\frac{\as}{2\pi}\Gamma O^R$）。则 $\langle j|O_i|j\rangle$ 的紫外极点给出 $\Gamma_{qq}=\gamma_{qq}(N)$，$\Gamma_{q_kg}=2\gamma_{qg}(N)$，$\Gamma_{gq}=\gamma_{gq}(N)$，$\Gamma_{gg}=\gamma_{gg}(N)$（偶数 $N$；奇数 $N$ 只有非单态 $\Gamma_{qq}=\gamma_{qq}(N)$）。Feynman 规范中逐图的值见\cref{tab:pdf-localgraphs}；单个图依赖于规范参数 $\xi_g$，只有包括 $\delta Z_2$、$\delta Z_3$ 的总和与 $\xi_g$ 无关。$N=2$ 时 $\Gamma$ 的两列之和为零：$T^{++}=\sum_kO^{(2)}_{q_k}+O^{(2)}_g$ 不被重整化。（本命题只涉及部分子态与算符。）
% VF-PDF-14, VF-PDF-10
\uses{def:dis-twist,thm:prelim-composite-renorm,prop:pdf-lo-kernels}
\end{proposition}

\begin{proof}
\emph{重整化条件。}树图 $\langle j|O_i|j\rangle^{(0)}=\delta_{ij}\cdot2(n_B\cdot p)^N$（偶数 $N$）。由\cref{thm:prelim-composite-renorm}，对物理外态扭度 2 算符只在 $\{O_q,O_g\}$ 内混合，于是 LSZ 振幅（含外线留数）的有限性要求 $\mathrm{UV}[\Gamma^{(1)}_{\mathrm{1PI}}+\delta Z_j\langle j|O_i|j\rangle^{(0)}]/(2(n_B\cdot p)^N)=\frac{\as}{2\pi}\frac1\varepsilon\Gamma_{ij}$，$\delta Z_q=\delta Z_2$，$\delta Z_g=\delta Z_3$。
\emph{紫外极点的提取。}单圈紫外发散定域、与红外正规化无关。对无质量在壳运动学，Feynman 参数化并平移后分母是 $q^2$ 的纯幂，只有分子中关于回路动量 $q$ 为二次（三个传播子）或零次（两个传播子）的项是对数发散的；用 $\langle q^\mu q^\nu\rangle=q^2g^{\mu\nu}/D$ 与 $\int\frac{\dd^Dq}{(2\pi)^D}\frac{q^2}{(q^2)^3}\big|_{\mathrm{UV}}=\int\frac{\dd^Dq}{(2\pi)^D}\frac{1}{(q^2)^2}\big|_{\mathrm{UV}}=\frac{\ii}{16\pi^2\varepsilon}$。下面记 $n=n_B$。
\emph{夸克部分。}顶点图 A1 $=-\ii g^2C_F\int\frac12\TrD[\sla p\gamma^\mu\sla k\sla n\sla k\gamma_\mu](n\cdot k)^{N-1}/((k^2)^2(p-k)^2)$；用 $\gamma^\mu\sla k\sla n\sla k\gamma_\mu=-(D-2)\sla k\sla n\sla k$、权重 $2(1-y)$、$k=q+yp$，分子的二次部分为 $2y^{N-1}(n\cdot p)^Nq^2$，极点为 $\frac{\as}{4\pi}\frac{C_F}\varepsilon(n\cdot p)^N\int_0^12(1-y)2y^{N-1}\dd y$，相对树图为 $\frac{\as}{2\pi}\frac{C_F}\varepsilon\frac1{N(N+1)}$。算符—胶子图 A2 $=$ A3 $=\ii g^2C_F\int\frac12\TrD[\sla p\sla n\sla k\sla n]\sum_{j=0}^{N-2}(n\cdot p)^j(n\cdot k)^{N-2-j}/(l^2k^2)$，$k=p-l$；由 $\sla n\sla k\sla n=2(n\cdot k)\sla n$ 与 $l=q+xp$，只剩与 $q$ 无关的项，极点为 $-\frac{\as}{4\pi}\frac{C_F}\varepsilon4(n\cdot p)^N\sum_{m=2}^N\frac1m$，即每个图 $-\frac{\as}{2\pi}\frac{C_F}{\varepsilon}(S_1(N)-1)$。加上 Feynman 规范的 $\delta Z_2=-\frac{\as}{4\pi}\frac{C_F}{\varepsilon}$：$\Gamma_{qq}=C_F[\frac{1}{N(N+1)}-2(S_1(N)-1)-\frac12]=\gamma_{qq}(N)$。
\emph{混合与胶子部分。}$\langle q|O_g|q\rangle$ 的单个图 B 给出 $\gamma_{gq}$；$\langle g|O_q|g\rangle$ 中夸克圈的两种走向 C1、C2 各给出 $\gamma_{qg}$，且由 C 宇称 $C_2=(-1)^NC_1$，故偶数 $N$ 时 $\Gamma_{q_kg}=2\gamma_{qg}$，奇数 $N$ 时为零（$N=1$ 即 Furry 定理）。D1（算符的双胶子顶点与两个三胶子顶点）、D2 $=$ D3（算符的三胶子顶点）与 $\delta Z_3=\frac{\as}{4\pi}\frac1\varepsilon[(\frac{13}6-\frac{\xi_g}2)C_A-\frac43T_Fn_f]$ 之和为 $\gamma_{gg}$；$n_f$ 部分全部来自 $\delta Z_3$。由算符双胶子项构成的蝌蚪图与“算符双胶子顶点 $+$ 四胶子顶点”图没有紫外极点。
\emph{规范无关性与 $N=2$。}表中 $\xi_g$ 依赖之和为零（夸克：$\frac{\xi_g-1}2-\frac{\xi_g-1}2=0$；胶子 $C_A$ 部分：$\frac{\xi_g-1}2-2\cdot\frac{\xi_g-1}8-\frac{\xi_g-1}4=0$）。$N=2$ 时 $\Gamma=\begin{pmatrix}-\frac43C_F&\frac23n_fT_F\\\frac43C_F&-\frac23n_fT_F\end{pmatrix}$（基 $\sum_kO_{q_k}$、$O_g$），两列之和为零。
% (graph values cross-checked for N=1..14 with explicit 4x4 Dirac matrices and with D-dim traces; writing/work/W09/v9_local_ops.py)
\end{proof}

\begin{table}[htbp]
\centering
\caption{扭度 2 定域算符单圈矩阵元的逐图紫外极点（单位 $\frac{\as}{2\pi}\frac1\varepsilon\times$ 色因子，相对于树图 $2(n_B\cdot p)^N$），以及一般协变规范中相对于 Feynman 规范（$\xi_g=1$）的改变。}\label{tab:pdf-localgraphs}
\small
\renewcommand{\arraystretch}{1.2}
\begin{tabular}{@{}lll@{}}
\toprule
图 & Feynman 规范 & 一般协变规范的改变\\
\midrule
A1 顶角（$C_F$） & $\frac{1}{N(N+1)}$ & $+\frac{\xi_g-1}{2}$\\
A2 $=$ A3 算符—胶子（$C_F$） & $-(S_1(N)-1)$ & 无\\
$\delta Z_2$（$C_F$） & $-\frac12$ & $-\frac{\xi_g-1}{2}$\\
B（$C_F$） & $\frac{N^2+N+2}{(N-1)N(N+1)}$ & 无\\
C1 $=$ C2（$T_F$，每味） & $\frac{N^2+N+2}{N(N+1)(N+2)}$ & 无（无内胶子线）\\
D1 双胶子算符顶角（$C_A$） & $\frac{2}{N(N-1)}+\frac{2}{(N+1)(N+2)}+\frac{1}{N+1}$ & $+\frac{\xi_g-1}{2}$\\
D2 $=$ D3 三胶子算符顶角（$C_A$） & $-(S_1(N)-1)-\frac{N+2}{2(N+1)}$ & $-\frac{\xi_g-1}{8}$\\
$\delta Z_3$ & $\frac56C_A-\frac23T_Fn_f$ & $-\frac{\xi_g-1}{4}C_A$\\
\bottomrule
\end{tabular}
\end{table}

\begin{example}[精确到 $\varepsilon$ 的单圈矩]\label{ex:pdf-moments-exact}
取\cref{eq:pdf-bare1}的 $n$ 阶矩，并与同样以胶子质量正规化的定域算符单圈矩阵元比较，PDF 图与定域算符图逐一对应且对 $\varepsilon$ 精确相等（单位 $A_\varepsilon$）：梯形图 R1 $\leftrightarrow$ 顶角图 A1，$(1-\varepsilon)^2\int_0^1\xi^{n-1}(1-\xi)\xi^{-\varepsilon}=\frac{(1-\varepsilon)^2}{(n-\varepsilon)(n+1-\varepsilon)}$；Wilson 线图 R2$+$V1 $\leftrightarrow$ 算符—胶子图 A2$+$A3，$\int_0^1\frac{2\xi^{1-\varepsilon}(\xi^{n-1}-1)}{1-\xi}=-2\sum_{i=1}^{n-1}\frac{1}{i+1-\varepsilon}$；自能 $\leftrightarrow$ 自能，$-\frac{1-\varepsilon}{2-\varepsilon}$。
% VF-PDF-13
总和：$n=1$ 时为 $0$（夸克数守恒对任意 $\varepsilon$ 成立）；$n=2$ 时 $\frac{(1-\varepsilon)^2-2(3-\varepsilon)-(1-\varepsilon)(3-\varepsilon)}{(2-\varepsilon)(3-\varepsilon)}=-\frac{4}{3-\varepsilon}$，极点 $-\frac43\frac{\as}{2\pi}\frac{C_F}\varepsilon=\frac{\as}{2\pi}\frac{\gamma_{qq}(2)}\varepsilon$；$n=3$ 时 $-\frac{2(5-2\varepsilon)^2}{(4-\varepsilon)(3-\varepsilon)(2-\varepsilon)}$，极点系数 $-\frac{25}{12}=\gamma_{qq}(3)/C_F$。一般地总和在 $\varepsilon\to0$ 的值 $\frac{1}{n(n+1)}-2(S_1(n)-1)-\frac12=\gamma_{qq}(n)/C_F$，于是 PDF 抵消项 $Z(x)=\delta(1-x)-\frac{\as}{2\pi}\frac1\varepsilon P_{qq}$ 的矩恰是定域算符的 $Z_O(n)$——\cref{thm:pdf-moments-local}第 (2) 步在单圈的显式实例。
\uses{thm:pdf-moments-local,prop:pdf-pqq-operator}
\end{example}

\begin{example}[胶子中的夸克]\label{ex:pdf-quark-in-gluon}
单圈 $f_{q/g}$ 来自 $g\to q\bar q$，$D$ 维核 $\hat P_{qg}=T_F[1-\frac{2z(1-z)}{1-\varepsilon}]$（\cref{prop:pdf-lo-kernels}）。无质量夸克的纯维数正规化中 $\MSbar$ 分布是纯红外极点 $-\frac{\as}{2\pi}\frac1\varepsilon P_{qg}$。为看清极化平均约定的作用，给夸克质量 $m$，$A^+=0$ 规范的直接计算给出横向被积函数 $\frac{g^2T_F}{(2\pi)^{D-1}}\frac{(\hat P_{qg}/T_F)\bm l_{\mathrm T}^2+m^2}{(\bm l_{\mathrm T}^2+m^2)^2}$（$m^2$ 项来自分子中的夸克质量），于是
\[
f^{(1)}_{Q/g}(z)=\frac{\as}{2\pi}\ee^{\varepsilon\gamma_E}\Big(\frac{\mu^2}{m^2}\Big)^{\varepsilon}\big[(1-\varepsilon)\hat P_{qg}\,\Gamma(\varepsilon)+T_F\Gamma(1+\varepsilon)\big]
=\frac{\as}{2\pi}P_{qg}(z)\,\ee^{\varepsilon\gamma_E}\Gamma(\varepsilon)\Big(\frac{\mu^2}{m^2}\Big)^{\varepsilon},
\]
因为 $(1-\varepsilon)\hat P_{qg}=P_{qg}-\varepsilon T_F$ 而 $\varepsilon\Gamma(\varepsilon)=\Gamma(1+\varepsilon)$：质量项恰好抵消了 $D$ 维核的 $\order{\varepsilon}$ 部分。$\MSbar$ 重整化后 $f^{(1)}_{Q/g}=\frac{\as}{2\pi}P_{qg}(z)\ln\frac{\mu^2}{m^2}$，没有常数项；若把胶子极化平均改为 $\frac12$，结果乘以 $\frac{D-2}{2}=1-\varepsilon$，$\MSbar$ 分布变为 $\frac{\as}{2\pi}P_{qg}[\ln\frac{\mu^2}{m^2}-1]$：极化平均约定是 $\MSbar$ 方案的一部分（\cref{def:pdf-gluon}）。$\mu<m$ 时这个 $\MSbar$ 重夸克分布是负的：$\MSbar$ 分布不必为正，微扰论中就有初等的例子（\cref{conj:pdf-positivity}）。
\uses{def:pdf-gluon,prop:pdf-lo-kernels}
\end{example}

% -----------------------------------------------------------------------------------------------------
\section{$x$ 空间的 DIS 定理（扭度 2 部分）}\label{sec:pdf-dis}

\cref{thm:dis-moments}是关于矩的陈述，\cref{thm:pdf-moments-local}把其中的 $A^{(n)}_i$ 认同为 PDF 的矩；加上矩问题的唯一性（\cref{app:appK}），就得到结构函数扭度 2 部分在 $x$ 空间的因子化；被控制的只是扭度 2 部分，余项只在矩的意义下是 $\order{1/Q^2}$。

\begin{theorem}[DIS 的 $x$ 空间因子化（扭度 2 部分）\tagR]\label{thm:pdf-dis-xspace}
（$R_{\mathrm{OPE}}$：以 OPE 为前提而严格。QCD 中含夸克的微扰 OPE 是全阶论证（\cref{thm:dis-ope}），严格证明只对标量理论与纯 Yang--Mills 的规范不变算符给出；对强子态还须非微扰 OPE（\cref{asm:intro-np-ope}）。）
定义 $F_2^{\mathrm{tw2}}(x,Q^2)$ 为偶数矩满足 $\int_0^1\dd x\,x^{n-2}F_2^{\mathrm{tw2}}=\sum_iC^{(n)}_{2,i}\big(\frac{Q^2}{\mu^2},\as(\mu)\big)A^{(n)}_i(\mu)$（$n=2,4,\dots$）的分布；它由这些矩唯一确定（对强子 PDF 所属的可积函数类由\cref{cor:appK-pdf}，对部分子靶所属的分布类见证明）。则
\begin{equation}\label{eq:pdf-xspace}
F_2^{\mathrm{tw2}}(x,Q^2)/x=\textstyle\sum_i\big(C_{2,i}\otimes f_i\big)(x),
\end{equation}
$C_{2,i}$ 是 Mellin 矩为 $C^{(n)}_{2,i}$ 的唯一分布，构造地等于 $\MSbar$ 减除了共线极点的部分子结构函数。在 $\order{\as}$，记 $L=\ln(Q^2/\mu^2)$，
\begin{equation}\label{eq:pdf-F2nlo}
\frac{F_2^{\mathrm{tw2}}}{x}=\sum_qe_q^2\Big\{(q+\bar q)+\frac{\as}{2\pi}\Big[(C_2^q+L\,P_{qq})\otimes(q+\bar q)+2\,(C_2^g+L\,P_{qg})\otimes g\Big]\Big\}+\order{\as^2},
\end{equation}
\begin{equation}\label{eq:pdf-C2}
\begin{aligned}
C_2^q&=C_F\Big[2\mathcal D_1-\tfrac32\mathcal D_0-(1+z)\ln(1-z)-\tfrac{1+z^2}{1-z}\ln z+3+2z-\big(\tfrac92+\tfrac{\pi^2}3\big)\delta(1-z)\Big],\\
C_2^g&=T_F\Big[\big(z^2+(1-z)^2\big)\ln\tfrac{1-z}{z}-8z^2+8z-1\Big],
\end{aligned}
\end{equation}
% VF-DIS-7
$\mathcal D_k=[\ln^k(1-z)/(1-z)]_+$；纵向结构函数为 $\frac{F_L^{\mathrm{tw2}}}{x}=\frac{\as}{2\pi}\sum_qe_q^2[C_L^q\otimes(q+\bar q)+2C_L^g\otimes g]+\order{\as^2}$，$C_L^q=2C_Fz$，$C_L^g=4T_Fz(1-z)$。$F_2-F_2^{\mathrm{tw2}}$ 的每个固定偶数矩为 $\order{1/Q^2}$（强子：$n\ge n_0$；在\cref{asm:intro-growth}之下一切偶数 $n\ge2$）。
非微扰地位：\cref{asm:intro-wightman}、\cref{asm:intro-composite}、\cref{asm:intro-np-ope}、\cref{asm:intro-af-coeff}、\cref{asm:intro-growth}（MA0、MA3、MA4）。
\uses{thm:dis-moments,thm:pdf-moments-local,lem:prelim-mellin,cor:appK-pdf,thm:appK-even-moments,prop:appK-kernels,prop:pdf-partonic-dr,asm:intro-wightman,asm:intro-composite,asm:intro-np-ope,asm:intro-af-coeff,asm:intro-growth}
\end{theorem}

\begin{proof}
(1) 由\cref{thm:dis-moments}，对偶数 $n$，$\int_0^1x^{n-2}F_2=\sum_iC^{(n)}_{2,i}A^{(n)}_i+\order{1/Q^2}$（部分子靶：一切偶数 $n\ge2$；强子：$n\ge n_0$，或在\cref{asm:intro-growth}之下一切偶数 $n$）。(2) 由\cref{thm:pdf-moments-local}，偶数 $n$ 时 $A^{(n)}_q=\int_0^1x^{n-1}(q+\bar q)$，$A^{(n)}_g=\int_0^1x^{n-1}g$。(3) \emph{$C_{2,i}$ 的构造。}把 (1)(2) 用于纯维数正规化中的无质量在壳部分子靶 $j$：除 $Q$ 外无别的标度，幂次修正为零；由\cref{prop:pdf-partonic-dr}，$f_{i/j}=Z_{ij}$，故部分子结构函数满足 $\int_0^1z^{n-2}\hat F_{2,j}=\sum_iC^{(n)}_{2,i}Z_{ij}(n)$。令 $C_{2,\cdot}\equiv(\hat F_{2,\cdot}/z)\otimes Z^{-1}$（味道矩阵），它支撑在 $[0,1]$ 上，由\cref{lem:prelim-mellin}其 Mellin 矩恰为 $C^{(n)}_{2,i}$。(4) \emph{唯一性。}每一阶的 $C_{2,i}$ 属于分布类 $\Delta=r+\sum_{k\le K}a_k\mathcal D_k+b\,\delta(1-z)$，$|r(z)|\le Cz^{-1}(1+|\ln z|)(1+|\ln(1-z)|)^K$。若 $\Delta$ 的偶数矩全为零：沿偶数 $N\to\infty$，$\int z^{N-1}r\to0$，而 $\int_0^1z^{N-1}\mathcal D_k=\frac{(-1)^{k+1}}{k+1}\ln^{k+1}\bar N+(\text{低次})$，$\bar N=N\ee^{\gamma_E}$，从最高次起逐个得 $a_k=0$，再得 $b=0$，最后由\cref{prop:appK-kernels}证明中的论证（\cref{thm:appK-even-moments}）得 $r=0$。对强子，$\sum_iC_{2,i}\otimes f_i$ 是满足 $x^{N_1-1}(\cdot)\in L^1$ 的函数，唯一性即\cref{cor:appK-pdf}。(5) 由\cref{lem:prelim-mellin}，$\int_0^1x^{n-1}(C_{2,i}\otimes f_i)=C^{(n)}_{2,i}\int_0^1x^{n-1}f_i$，故\cref{eq:pdf-xspace}右边的偶数矩为 $\sum_iC^{(n)}_{2,i}A^{(n)}_i$，由 (4) 它就是 $F_2^{\mathrm{tw2}}/x$。(6) \emph{$\order{\as}$。}部分子结构函数为 $\hat F_2^q/(ze_q^2)=\delta(1-z)+\frac{\as}{2\pi}(\frac{\mu^2}{Q^2})^\varepsilon[-\frac1\varepsilon P_{qq}+C_2^q+\order\varepsilon]$、$\hat F_2^g/(ze_q^2)=\frac{\as}{2\pi}(\frac{\mu^2}{Q^2})^\varepsilon[-\frac2\varepsilon P_{qg}+2C_2^g+\order\varepsilon]$（胶子道的因子 $2$ 来自每味的夸克与反夸克）。按 (3) 乘以 $Z^{-1}=\delta+\frac{\as}{2\pi}\frac1\varepsilon P+\order{\as^2}$：$(1-\varepsilon L)(-\frac1\varepsilon P+C)+\frac1\varepsilon P=C+LP+\order\varepsilon$，即\cref{eq:pdf-F2nlo}。$F_L$ 由 $F_1$ 的矩同样得到，$C_L=C_2-C_1$ 在这一阶没有极点，与 $\mu$ 无关。(7) $F_2-F_2^{\mathrm{tw2}}$ 的偶数矩为 $\order{1/Q^2}$ 直接来自 (1)。
\end{proof}

两个检验：$C_2^q$、$C_2^g$、$C_L$ 的 Mellin 矩重现第~7~章单圈 Wilson 系数的闭式（如 $\int_0^1C_2^q=0$，即 Adler 求和规则；$C^q_{2,2}=\frac16C_F$，$C^g_{2,2}=-\frac14T_F$）；由\cref{thm:pdf-dglap}可验证 $\partial F_2^{\mathrm{tw2}}/\partial\ln\mu^2=\order{\as^2}$（与 $\partial L/\partial\ln\mu^2=-1$ 的项抵消）。定理控制的只是扭度 2 部分与余项的\emph{矩}。余项的逐点大小不受控制：$\sin(Q^2x)$ 的每个矩乘以 $Q^2$ 都有界而 $\sup|\sin(Q^2x)|=1$（\cref{rem:dis-nonuniform}）；OPE 所能给出的最好结果是弱收敛（\cref{cor:dis-weak-xspace}），而“存在矩为给定扭度 2 矩的测度”本身是 Hausdorff 矩问题（\cref{rem:appK-counterexamples}）。在固定 $0<x<1$ 处带 $\order{\Lambda^2/Q^2}$ 余项的逐点 $x$ 空间因子化是第~10~章的全阶论证（\cref{thm:disall-factorization}，\tagA），不是本定理的推论。

% -----------------------------------------------------------------------------------------------------
\section{求和规则、方案与正定性}\label{sec:pdf-sumrules}

\begin{proposition}[求和规则 \tagR]\label{prop:pdf-sum-rules}
（$R_{\mathrm{PT}}$：在重整化微扰论中对红外正规化的部分子态严格。）
(a) 数求和规则：$\int_0^1\dd x\,[q(x;\mu)-\bar q(x;\mu)]=N_q$ 对一切 $\mu$ 成立，$N_q$ 是靶的味道 $q$ 的净夸克数（守恒荷的本征值；对部分子靶 $j$，$N_q=\delta_{qj}-\delta_{\bar qj}$）；等价地 $\gamma_{qq}(1)=0$ 在各阶成立。
(b) 动量求和规则：$\sum_i\int_0^1\dd x\,x\,f_i(x;\mu)=1$ 在 $\MSbar$ 中对一切 $\mu$ 精确成立（$T^{++}$ 有限；$N=2$ 反常维度矩阵有本征值 $0$）。
非微扰地位：对强子需\cref{asm:intro-wightman}、\cref{asm:intro-composite}（MA0）。
\uses{thm:pdf-moments-local,prop:pdf-invariance,prop:pdf-lo-kernels,thm:pdf-dglap,asm:intro-wightman,asm:intro-composite}
\end{proposition}

\begin{proof}
(a) \cref{thm:pdf-moments-local}取 $n=1$：$\int_0^1(q-\bar q)=\int_{-1}^1f_q\,\dd\xi=\langle P|[\bar\psi_q\gamma^+\psi_q]_\mu|P\rangle/(2P^+)$。$j^\mu_q=\bar\psi_q\gamma^\mu\psi_q$ 是味道 $q$ 的数守恒流，其荷是好量子数、与标度无关，故该算符不需重整化（单圈即 $\gamma_{qq}(1)=0$，\cref{prop:pdf-lo-kernels}）。由平移不变性，$\int\dd x^-\dd^2\bm x_{\mathrm T}\langle P'|j^+_q(x)|P\rangle=(2\pi)^3\delta^{(3)}(P'-P)\langle P|j^+_q(0)|P\rangle$，而左边等于荷的本征值乘以 $\langle P'|P\rangle$，故 $\langle P|j^+_q(0)|P\rangle=2P^+N_q$。
(b) \cref{thm:pdf-moments-local}取 $n=2$ 并对各味求和：$\sum_q\int_0^1x(q+\bar q)+\int_0^1xg=[\sum_q\langle O^{(2)}_q\rangle+\langle O^{(2)}_g\rangle]_\mu/(2(P^+)^2)$。规范不变对称能动张量 $T^{\mu\nu}=\sum_q\frac{\ii}4\bar\psi_q(\gamma^\mu\overleftrightarrow D^\nu+\gamma^\nu\overleftrightarrow D^\mu)\psi_q-G^{\mu\alpha}_aG^\nu{}_{\alpha a}+\frac14g^{\mu\nu}G^2$（规范固定项与鬼场项是 BRST 恰当的，在物理态之间不贡献）的 $++$ 分量在前向矩阵元中等于 $\sum_qO^{(2)}_q+O^{(2)}_g$（夸克部分把导数移到右边，胶子部分见\cref{prop:pdf-invariance}(d)）。能动张量守恒，其荷为四动量，故有限、不需重整化（\cite{Collins1984}），且同 (a) 一样 $\langle P|T^{\mu\nu}|P\rangle=2P^\mu P^\nu$，于是左边等于 $\langle P|T^{++}|P\rangle/(2(P^+)^2)=1$。用 DGLAP 的语言，这等价于 $N=2$ 的全阶反常维度矩阵各列之和为零；单圈时即\cref{prop:pdf-lo-kernels}中的动量求和规则（\cref{thm:pdf-dglap}）。
\end{proof}

对强子，$N_q$ 是价夸克数（质子 $N_u=2$，$N_d=1$）；强子\emph{结构函数}的动量求和规则另需增长条件\cref{asm:intro-growth}，它与低阶矩的色散表示等价（\cref{prop:dis-low-moments}）。

\begin{remark}[解析延拓与 $P^\pm_{\mathrm{NS}}$]\label{rem:pdf-ns-plusminus}
定域算符只给出整数 $N$ 处的反常维度，电磁 DIS 中只出现偶数 $N$；$x$ 空间的核是某个解析延拓的 Mellin 逆变换。由偶数 $N$ 出发的延拓（C 偶，$q+\bar q$）与由奇数 $N$ 出发的延拓（C 奇，$q-\bar q$）一般是\emph{不同}的函数；在多项式有界的条件下各自唯一（\cref{prop:appK-continuation}），例如 $S_1(N)\mapsto\psi(N+1)+\gamma_E$，而加上 $c\sin(\pi N/2)$ 因指数增长被排除。领头阶 $qq$ 道两种延拓相同；从次领头阶起 $P^+_{\mathrm{NS}}-P^-_{\mathrm{NS}}=2P^V_{q\bar q}\ne0$，色因子为 $C_F(C_F-\frac{C_A}2)$\cite{CurciFurmanskiPetronzio1980}。以 $(\as/2\pi)^2$ 为展开参数，$\int_0^1P^{V(1)}_{q\bar q}\,\dd x=0.71931\,C_F(C_F-\frac{C_A}2)\ne0$，而 $\int_0^1P^-_{\mathrm{NS}}=0$（价夸克数守恒），所以 $P^+_{\mathrm{NS}}$ 的第一矩不为零。“DGLAP 方程与 OPE 重整化群方程在所有阶以单一的核等价”是错的：两者等价，但非单态有两个核，分别对应两种 C 宇称。
% VF-REF-5
\end{remark}

\begin{remark}[方案变换]\label{rem:pdf-schemes}
$\MSbar$ 只是约定。对有限变换 $f'=Z\otimes f$（$Z=\mathbb 1+\frac{\as}{2\pi}Z_1+\cdots$，与靶无关），系数函数变为 $C'=C\otimes Z^{-1}$，可观测量不变，演化核变为 $\mathcal P'=Z\mathcal PZ^{-1}+\frac{\dd Z}{\dd\ln\mu^2}Z^{-1}$，即 $P'^{(1)}=P^{(1)}+Z_1\otimes P^{(0)}-P^{(0)}\otimes Z_1-\frac{\beta_0}2Z_1$（$\mathcal P=\frac{\as}{2\pi}P^{(0)}+(\frac{\as}{2\pi})^2P^{(1)}$）。DIS 方案选 $Z_1$ 使 $F_2$ 保持部分子模型的形式：$q^{\mathrm{DIS}}_i=q_i+\frac{\as}{2\pi}[C_2^q\otimes q_i+C_2^g\otimes g]$，并按动量守恒令 $g^{\mathrm{DIS}}=g-\frac{\as}{2\pi}[C_2^q\otimes\Sigma+2n_fC_2^g\otimes g]$（参见\cite{DiemozEtAl1988}），于是 $P^{(1),\mathrm{DIS}}_{\mathrm{NS}}=P^{(1),\MSbar}_{\mathrm{NS}}-\frac{\beta_0}2C_2^q$，$C_L$ 在 $\order{\as}$ 与方案无关。胶子极化的平均约定（\cref{def:pdf-gluon}）也是方案的一部分。DIS 方案是合法的方案选择；但若把 PDF \emph{定义}为“$F_2$ 中的那个量”，其他过程中的普适性就成了定义而非定理。
% VF-DIS-11
\end{remark}

\begin{conjecture}[$\MSbar$ PDF 的正定性 \tagC]\label{conj:pdf-positivity}
存在一个（非微扰的）标度 $\mu_0$，使得对一切 $\mu\ge\mu_0$，强子的 $\MSbar$ 夸克、反夸克与胶子分布非负：$q(x;\mu),\,\bar q(x;\mu),\,g(x;\mu)\ge0$（$0<x<1$）。
\uses{def:pdf-quark,def:pdf-gluon}
\end{conjecture}

\paragraph{地位与缺口}
没有一般的证明，而更强的命题是错的。\emph{支持：}Candido、Forte 与 Hekhorn\cite{CandidoForteHekhorn2020}论证次领头阶 $\MSbar$ PDF 在微扰区域为正；Candido、Forte、Giani 与 Hekhorn\cite{Candido:2023ujx}同意 PDF 在足够低的非微扰标度会变负，并保守地半定量估计 $Q^2\gtrsim5\,\mathrm{GeV}^2$ 以上正定性得以继承。\emph{反对一般定理：}Collins、Rogers 与 Sato\cite[第~VI~节]{CollinsRogersSato2022}指出上述论证在“裸 PDF 为正”这一步失效——用维数正规化同时处理紫外（$n<4$）与共线（$n>4$）发散时，在 $n<4$ 为正的紫外贡献延拓到 $n>4$ 后可以变负。\emph{“在所有标度上非负”是错的：}Yukawa 模型中单圈 $\MSbar$ 分布为 $f_{q/p}(\xi;\mu)=a_\lambda(1-\xi)[\ln\frac{\mu^2}{\Delta(\xi)}-1+\frac{(m_q+\xi m_p)^2}{\Delta(\xi)}]$，$\Delta(\xi)=\xi m_s^2+(1-\xi)m_q^2-\xi(1-\xi)m_p^2$，$a_\lambda=\frac{\lambda^2}{16\pi^2}$\cite[式~(55)]{CollinsRogersSato2022}；
% VF-FRO-10
对 $m_q,m_p,m_s=0.3,1.0,1.5\,\mathrm{GeV}$，$\mu=0.8\,\mathrm{GeV}$ 时恰在 $\xi\in(0.264,1)$ 上 $f<0$，$\mu=1.2\,\mathrm{GeV}$ 时在 $\xi>0.567$ 上 $f<0$，而 $f$ 对一切 $\xi$ 为正当且仅当 $\mu\ge\mu_*=m_s\exp\{\frac12[1-(m_q+m_p)^2/m_s^2]\}=1.699\,\mathrm{GeV}$（这些数值是我们对该式的直接计算）。QCD 中也有初等的例子：单圈匹配给出 $\MSbar$ 重夸克分布 $f_h(x;\mu)=\frac{\as}{2\pi}\ln\frac{\mu^2}{m_h^2}(P_{qg}\otimes g)(x)+\order{\as^2}$（\cref{ex:pdf-quark-in-gluon}），$\mu<m_h$ 且 $g>0$ 时为负。\emph{循环论证的警示：}领头阶 DGLAP 演化保持正定性，并不证明初始条件为正，而高阶的正定性保持依赖于方案。因此猜想只针对 $\mu\ge\mu_0$，而 $\mu_0$ 是非微扰的；本书只陈述这些事实，不作进一步的裁决，详细讨论（含 Yukawa 曲线）见第~24~章。强子 $\MSbar$ PDF 本身的存在依赖于\cref{asm:intro-wightman}与\cref{asm:intro-composite}；此缺口与幂次计数无关，不涉及 E-S 至 E-W 型的误差估计。升级途径：一个证明或一个 $\mu\ge\mu_0$ 区域的反例，以及 $\mu_0$ 的非微扰确定。

% -----------------------------------------------------------------------------------------------------
\section{碎裂函数、类时演化与极化分布的定义}\label{sec:pdf-ff}

碎裂函数描述高能部分子产生动量份额为 $z$ 的强子 $h$ 的“数密度”。其算符定义与 PDF 平行，但被观测的强子在\emph{末态}，对其余末态 $X$ 的求和夹在两个场之间，并由此带来参考系与归一问题。

\begin{definition}[碎裂函数]\label{def:pdf-ff}
设夸克动量 $k$ 沿 $+z$ 方向，碎裂出强子 $h$，$z=p_h^+/k^+$。在强子系（$\bm p_{h,\mathrm T}=0$）中取 $w=(0^+,w^-,\bm 0_{\mathrm T})$，
\begin{equation}\label{eq:pdf-ff}
\begin{aligned}
d_{h/q}(z)=\frac{z^{1-2\varepsilon}}{4N_c}\sum_X\int\frac{\dd w^-}{2\pi}\,\ee^{\ii k^+w^-}\Trc\TrD\Big[&\gamma^+\langle0|W[\infty n_B,w]\psi(w)|hX\rangle\\
&\times\langle hX|\bar\psi(0)W[0,\infty n_B]|0\rangle\Big],
\end{aligned}
\end{equation}
% VF-CONV-17
Wilson 线是沿 $n_B$ 的未来指向线。等价地，在部分子系（$\bm k_{\mathrm T}=0$）中前置因子为 $\frac1{4N_cz}$，并显式对 $\int\dd^{2-2\varepsilon}\bm p_{h,\mathrm T}$ 积分；两者由 $\dd^{2-2\varepsilon}\bm p_{h,\mathrm T}=z^{2-2\varepsilon}\dd^{2-2\varepsilon}\bm k_{\mathrm T}$ 联系。胶子碎裂函数类似地定义为
\begin{equation}\label{eq:pdf-ffg}
\begin{aligned}
d_{h/g}(z)=\frac{z^{1-2\varepsilon}}{(D-2)(N_c^2-1)k^+}\sum_X\int\frac{\dd w^-}{2\pi}\ee^{\ii k^+w^-}\sum_{j,a}&\langle0|W^{(A)}_{ab}[\infty n_B,w]G^{+j}_b(w)|hX\rangle\\
&\times\langle hX|G^{+j}_c(0)W^{(A)}_{ca}[0,\infty n_B]|0\rangle,
\end{aligned}
\end{equation}
$W^{(A)}$ 如\cref{def:pdf-gluon}。自由部分子给出 $d^{(0)}_{q/q}=d^{(0)}_{g/g}=\delta(1-z)$。碎裂函数在 Mellin 空间中按 $\MSbar$ 重整化。
\uses{def:prelim-wilson,def:prelim-lightfront,def:pdf-gluon}
\end{definition}

\paragraph{归一的检验。}
强子系由部分子系经横向推动得到，部分子的横动量为 $\bm k_{\mathrm T}=-\bm p_{h,\mathrm T}/z$，故 $\frac1z\cdot z^{2-2\varepsilon}=z^{1-2\varepsilon}=z^{d-3}$，即\cite[第~12.4~节，式~(12.38)、(12.40)]{Collins2011}中积分碎裂函数所带的 $z^{n-3}$。自由夸克检验（$\Trc\TrD[\gamma^+u\bar u]=4N_cp_h^+$，$\int\frac{\dd w^-}{2\pi}\ee^{\ii(k^+-p_h^+)w^-}=\delta(1-z)/k^+$）只在 $z=1$ 处取值，看不出 $z$ 的幂次；能看出的是 $\order{\as}$ 的检验：$A^+=0$ 规范中取 $h=q$、$X=g$，振幅 $\langle0|\psi(0)|q(p_h)g(l)\rangle=-g\,t^a\frac{\sla k\sla\epsilon}{k^2}u(p_h)$，$k=p_h+l$，强子系中 $p_h=(zk^+,0^-,\bm 0_{\mathrm T})$，$k^2=z\bm l_{\mathrm T}^2/(1-z)$，迹与相空间给出
\[
d^{(1)}_{q/q}(z)=\frac{g^2}{(2\pi)^{D-1}}\int\dd^{D-2}\bm l_{\mathrm T}\frac{z^{-2\varepsilon}\hat P_{qq}(z;\varepsilon)}{\bm l_{\mathrm T}^2}\ \xrightarrow{\ D=4\ }\ \frac{\as}{2\pi}C_F\frac{1+z^2}{1-z}\int\frac{\dd\bm l_{\mathrm T}^2}{\bm l_{\mathrm T}^2}\quad(z<1),
\]
恰是领头阶的类时分裂概率；若在强子系中误用 $\frac1{4N_cz}$，结果会多出 $z^{-2+2\varepsilon}$。对纯维数正规化中的无质量部分子，裸部分子碎裂函数无论前置因子如何都是 $\delta(1-z)$（无标度），故 $\MSbar$ 部分子系数函数不受影响；$z^{-2\varepsilon}$ 固定的是强子层次定义的方案。

\begin{remark}[碎裂函数的乘法可重整性是独立输入]\label{rem:pdf-ff-multren}
碎裂函数的 Mellin 矩\emph{不是}定域算符的矩阵元：\cref{eq:pdf-ff}中 $\sum_X|hX\rangle\langle hX|$ 夹在两个场之间，矩不能化为 $w^-=0$ 处的导数，\cref{thm:pdf-moments-local}的论证不适用。因此碎裂函数在矩空间中的乘法可重整性是独立的输入，即\cref{asm:appN-multren}(Z1) 对碎裂函数的情形；类时 DGLAP 在已计算的阶上是 $[R,R_{\mathrm{PT}}]$，作为全阶算符陈述则是 \tagA（\cref{prop:pdf-timelike-dglap}）。动量求和规则 $\sum_h\int_0^1\dd z\,z\,d_{h/q}=1$ 可由算符定义作形式推导（Collins 与 Soper 的原文讨论了这类求和规则，据该文摘要\cite{CollinsSoper1982a}），但推导要交换对强子的求和与重整化，本书视之为形式关系；据 Collins 与 Rogers 的摘要\cite{Collins:2023cuo}，夸克碎裂函数的荷（数）型求和规则在 QCD 中可能不成立，原因是 Wilson 线导致位于 $z=0$ 的束缚态，而这一效应不影响动量求和规则。
\uses{def:pdf-ff,thm:pdf-moments-local}
\end{remark}

\begin{proposition}[类时 DGLAP 演化 \tagA]\label{prop:pdf-timelike-dglap}
(a) 在碎裂函数的乘法可重整性（\cref{asm:appN-multren}(Z1)）之下，对一切阶
\begin{equation}\label{eq:pdf-tdglap}
\frac{\partial d_{h/i}(z;\mu)}{\partial\ln\mu^2}=\sum_j\big(P^T_{ji}\otimes d_{h/j}\big)(z),\qquad P^T_{ji}=\frac{\as}{2\pi}P^{T(0)}_{ji}+\Big(\frac{\as}{2\pi}\Big)^2P^{T(1)}_{ji}+\cdots,
\end{equation}
指标结构与 PDF 的演化互为转置；$P^{T(1)}\ne P^{(1)}$\cite{CurciFurmanskiPetronzio1980}；动量求和规则是形式的。
(b)（单圈内容，\tagR）（$R_{\mathrm{PT}}$：在重整化微扰论中对红外正规化的部分子态严格。）在单圈阶，部分子碎裂函数 $d_{j/i}$ 的紫外极点给出 $P^{T(0)}_{ji}=P_{ji}$（Gribov--Lipatov 关系\cite{GribovLipatov1972b}）；在 Feynman 规范与胶子质量正规化下，单圈裸碎裂函数与单圈裸 PDF 逐图相等：$d^{(1)}_{q/q}(z)=f^{(1)}_{q/q}(z)$。
非微扰地位：强子态需\cref{asm:intro-wightman}、\cref{asm:intro-composite}（MA0）。
\uses{def:pdf-ff,rem:pdf-ff-multren,prop:pdf-lo-kernels,asm:appN-multren,asm:intro-wightman,asm:intro-composite}
\end{proposition}

\begin{proofsketch}[(a) 的证明概要]
(1) 由\cref{asm:appN-multren}(Z1)，矩空间中 $d^{\mathrm{ren}}_{h/i}(N)=\sum_jd^{\mathrm{bare}}_{h/j}(N)Z^T_{ji}(N)$，$Z^T(N)$ 是与 $h$ 无关的纯极点矩阵。(2) 裸碎裂函数与 $\mu$ 无关；与\cref{thm:pdf-dglap}的证明相同的代数给出 $\partial d_{h/i}(N)/\partial\ln\mu^2=\sum_jd_{h/j}(N)P^T_{ji}(N)$，一致性条件保证 $P^T$ 有限且 $P^T=-\as\partial Z^{T(1)}/\partial\as$。(3) $P^T(N)$ 在 $\Real N>\sigma_0$ 中多项式有界，由偶数（或奇数）$N$ 的值唯一确定，Mellin 逆变换给出 $x$ 空间的卷积形式。(4) 单圈由 (b) 得 $P^{T(0)}=P^{(0)}$；从次领头阶起 $P^T\ne P$。碎裂函数的定义见\cite{CollinsSoper1982a}与\cite[第~12.4~节]{Collins2011}。
\end{proofsketch}

\begin{proof}[(b) 的证明]
(i) 夸克道，Feynman 规范，胶子质量 $m$，强子系：$p_h=(zk^+,0^-,\bm 0_{\mathrm T})$，$l^2=m^2$，$k=p_h+l$，$k^2=(z\bm l_{\mathrm T}^2+m^2)/(1-z)$。实图与\cref{fig:pdf-oneloop}一一对应（割线隔开被观测的夸克与胶子）；振幅 $\langle0|W[\infty n_B,0]\psi(0)|q(p_h)g(l)\rangle=g\,t^a\epsilon_\mu[-\frac{\sla k\gamma^\mu}{k^2}+\frac{n_B^\mu}{n_B\cdot l}]u(p_h)$（$\epsilon\to l$ 时 $-\sla k\sla lu/k^2+u=0$）。乘以 $z^{1-2\varepsilon}/(4N_c)$ 与 $1/(2l^+)$，横向被积函数（单位 $\frac{g^2C_F}{(2\pi)^{D-1}}\dd^{D-2}\bm l_{\mathrm T}$）为：梯形图 $\frac{z^{-2\varepsilon}(1-\varepsilon)(1-z)\bm l_{\mathrm T}^2}{(\bm l_{\mathrm T}^2+m^2/z)^2}$，Wilson 线图（含镜像）$\frac{2z^{1-2\varepsilon}}{(1-z)(\bm l_{\mathrm T}^2+m^2/z)}$；以 $M^2=m^2/z$ 积分得 $A_\varepsilon(1-\varepsilon)^2(1-z)z^{-\varepsilon}$ 与 $A_\varepsilon\frac{2z^{1-\varepsilon}}{1-z}$，与\cref{prop:pdf-pqq-operator}(a) 在 $\xi=z$ 处逐项相同。虚图不穿过割线，是同一振幅 $\langle0|[W[\infty n_B,0]\psi](0)|q(p)\rangle$ 的单圈修正（单粒子态的入态与出态相同），与 PDF 的虚图一样。因此 $d^{(1)}_{q/q}=f^{(1)}_{q/q}$ 对 $\varepsilon$ 精确成立，紫外极点为 $\frac{\as}{2\pi}\frac1\varepsilon P_{qq}$，包括 $\frac32C_F\delta(1-z)$。
(ii) 其余道。$A^+=0$ 规范中强子系的单圈实辐射密度对四个道都满足 $\bm l_{\mathrm T}^2\times(\text{密度})=z^{-2\varepsilon}\hat P_{ji}(z;\varepsilon)$（$q\to g$ 用夸克碎裂函数并对被观测胶子的物理极化求和，$g\to q$、$g\to g$ 用\cref{eq:pdf-ffg}），而 PDF 的相应密度是 $\hat P_{ji}(z;\varepsilon)$；两者的 $1/\varepsilon$ 系数逐点相同，故 $z<1$ 时 $P^{T(0)}_{ji}=P_{ji}$。$gg$ 道的 $\delta(1-z)$ 项：虚图仍是同一振幅 $\langle0|[W^{(A)}G^{+j}](0)|g(p)\rangle$ 的修正，与 $f_{g/g}$ 相同；实图之差 $(z^{-2\varepsilon}-1)\times(\text{PDF 实图})$ 乘以 $1/\varepsilon$ 后有限，且因 $(z^{-2\varepsilon}-1)/(1-z)\to2\varepsilon$ 在 $z\to1$ 可积，不改变作为分布的紫外极点，故 $P^{T(0)}_{gg}=P_{gg}$，包括 $\frac{\beta_0}2\delta(1-z)$。结合\cref{prop:pdf-lo-kernels}即得 (b)。
\end{proof}

\paragraph{地位与缺口}
全阶陈述 (a) 依赖于\cref{asm:appN-multren}(Z1) 对碎裂函数的情形（\tagA 级）：碎裂函数的矩不是定域算符的矩阵元，\cref{sec:pdf-moments}的论证不适用（\cref{rem:pdf-ff-multren}），这一步因此作为编号假设列出，而非从定域算符的重整化推出。此缺口与幂次计数无关，不涉及 E-S 至 E-W 型的误差估计。(b) 是严格的单圈结果。若能证明碎裂算符在矩空间中的乘法可重整性（例如一个类时算符的定义与相应的全阶图形论证），(a) 即可不再需要这一技术假设。Gribov--Lipatov 关系只在领头阶成立：从次领头阶起 $\MSbar$ 的类时与类空分裂函数不同\cite{CurciFurmanskiPetronzio1980}。

\begin{definition}[螺旋度与横度分布]\label{def:pdf-polarized}
对螺旋度 $+\frac12$ 的纵向极化靶，夸克螺旋度分布为
\begin{equation}\label{eq:pdf-helicity}
\Delta f_{q/H}(\xi)=\int\frac{\dd w^-}{4\pi}\ee^{-\ii\xi P^+w^-}\langle P,+|\bar\psi(w)\gamma^+\gamma_5W[w,0]\psi(0)|P,+\rangle_{\mathrm c},
\end{equation}
$x>0$ 时 $\Delta q(x)=\Delta f_q(x)$，$\Delta\bar q(x)=+\Delta f_q(-x)$；$\gamma_5=\ii\gamma^0\gamma^1\gamma^2\gamma^3$（\cref{def:prelim-gamma5}）；对螺旋度 $h$ 的自由夸克 $\bar u_h\gamma^+\gamma_5u_h=(2h)2p^+$，故 $\Delta q=q^\uparrow-q^\downarrow$；$\bar\psi\gamma^+\gamma_5\psi$ 是 C 偶的，所以反夸克关系的符号与\cref{def:pdf-antiquark}相反。对横向极化靶（横向自旋矢量 $\bm S_{\mathrm T}$），横度分布由
\begin{equation}\label{eq:pdf-transversity}
\int\frac{\dd w^-}{4\pi}\ee^{-\ii\xi P^+w^-}\langle P,S_{\mathrm T}|\bar\psi(w)\,\ii\sigma^{j+}\gamma_5W[w,0]\psi(0)|P,S_{\mathrm T}\rangle_{\mathrm c}=S_{\mathrm T}^j\,h_{1,q/H}(\xi)
\end{equation}
定义，$\ii\sigma^{j+}\gamma_5=\gamma^+\gamma^j\gamma_5$（$j$ 横向）；自旋沿 $+x$ 的自由夸克给出 $h_1=\delta(1-\xi)$。胶子螺旋度分布为
\begin{equation}\label{eq:pdf-dg}
\Delta f_{g/H}(\xi)=-\sum_{i,j}\int\frac{\dd w^-}{2\pi\xi P^+}\ee^{-\ii\xi P^+w^-}\,\ii\epsilon_{\mathrm T}^{ij}\langle P,+|G^{+i}_a(w)W^{(A)}_{ab}[w,0]G^{+j}_b(0)|P,+\rangle_{\mathrm c},
\end{equation}
其中 $\epsilon_{\mathrm T}^{12}=+1$；螺旋度 $+1$ 的自由胶子给出 $\Delta f_{g/g}=\delta(1-\xi)+\delta(1+\xi)$（偶分布，$\Delta g(x)=\Delta f_g(x)$）。维数正规化中 $\gamma_5$ 与 $\epsilon$ 张量按 HVBM 方案处理，并对轴矢流作 Larin 的有限重整化\cite{Larin:1993tq}，极化分布的 $\MSbar$ 方案由此确定。
\uses{def:pdf-quark,def:pdf-antiquark,def:pdf-gluon,def:prelim-gamma5}
\end{definition}

横度分布是手征奇的，单举 DIS 中不出现，需要另一个手征奇的对象才能测量（如横向极化 Drell--Yan\cite{RalstonSoper1979,JaffeJi1991,JaffeJi1992}）；自旋 $\frac12$ 的靶没有胶子横度分布。极化分布的演化（次领头阶核见\cite{MertigVanNeerven1996,Vogelsang1996}）与因子化见第~11~章。

% -----------------------------------------------------------------------------------------------------
\begin{tcolorbox}[enhanced,breakable,title=本章澄清的常见误解,fonttitle=\bfseries,
  colback=gray!3!white,colframe=gray!55!black,boxrule=0.6pt,arc=2pt,
  left=5pt,right=5pt,top=3pt,bottom=3pt,before skip=8pt,after skip=6pt]
\begin{enumerate}[leftmargin=1.6em,itemsep=1pt,topsep=1pt,label=\arabic*.]
\item 「PDF 由 DIS 定义。」——PDF 由算符定义；用 $F_2$ 定义会把普适性变成定义（\cref{rem:pdf-schemes}）。
\item 「PDF 是概率密度，故 $\MSbar$ PDF 为正」或「领头阶演化保持正定性，故 $\MSbar$ PDF 为正」。——数密度只对裸的截断量成立（\cref{rem:pdf-number-density}）；后一说法是循环论证，初始条件的正定性正是问题所在；$\MSbar$ 正定性是猜想，在所有标度上非负是错的（\cref{conj:pdf-positivity}）。
\item 「反常维度唯一确定 $x$ 空间的核。」——要选定解析延拓；偶数与奇数 $N$ 的延拓不同，$P^+_{\mathrm{NS}}\ne P^-_{\mathrm{NS}}$ 从次领头阶起（\cref{rem:pdf-ns-plusminus}）。
\item 「光锥规范去掉了所有 Wilson 线。」——剩余规范自由、$1/k^+$ 处方与无穷远横向连接仍在（\cref{rem:pdf-gauge-residual}）。
\item 「类时与类空分裂函数相同。」——只在领头阶成立（\cref{prop:pdf-timelike-dglap}）。
\item 「维数正规化中 $f_{q/q}=\delta(1-x)$，什么也没发生。」——裸分布是；重整化后它是纯红外极点（\cref{prop:pdf-partonic-dr}）。
\item 「DGLAP 来自截面的标度无关性。」——DGLAP 来自算符重整化（\cref{thm:pdf-dglap}），截面的标度无关性是因子化的推论。
\item 「PDF 的支撑是 $0\le\xi\le1$。」——$f_q$ 支撑在 $[-1,1]$，$\xi<0$ 部分是反号的反夸克密度（\cref{thm:pdf-support-pt}）。
\end{enumerate}
\end{tcolorbox}

\paragraph{本章小结。}
PDF 与碎裂函数在任何因子化定理之前就有了精确的算符定义：PDF 的矩是扭度 2 定域算符的矩阵元（\cref{app:appL}），DGLAP 由算符重整化得出，求和规则来自守恒流，DIS 扭度 2 部分的 $x$ 空间因子化以 OPE 为前提而严格；碎裂函数的全阶演化需要独立的乘法可重整性输入，$\MSbar$ PDF 的正定性只是猜想。第~9~章起，这些分布作为因子化定理中的长距离因子出现。
